\documentclass[dvipsnames]{ltjsarticle}
\usepackage{../../myStyle}
%% バッチ処理をするので\lheadは必ず4行目にあること
\lhead{version 3.0.18}
\begin{document}

\title{心理学データ解析(心理学統計法)}
\author{担当：小杉考司}
\date{Last Compiled on \number\year .\number\month .\number\day}
\maketitle

\setcounter{tocdepth}{1}
\newpage
\section*{はじめに}
\subsection*{授業のテーマ}

「見えないもの」をみよう・測ろう・考えようとするのが心理統計の世界である。

\subsection*{一年を通じて伝えたいポイント}
\begin{description}
	\item[母数を知るための推測] 大学での統計，心理統計は標本の記述統計量を検証するのではなく，母集団を表す数字を求めるための推測統計である。
	\item[確率を用いた推論] 母数を知るための方法は3つある。代表値を用いたモーメント法，確率分布をつかった最尤法，ベイズ法である。
	\item[要因計画と線形モデル] 心理学では，要因計画による研究条件をコントロールし，平均値の差をもちいた考察(平均因果効果)をおこなうことが多い。そしてそれらは総じて線形モデルとして表現できる。
	\item[モデル比較と意思決定] モデルによる母数の推定だけではなく，そこから一定の「結論」あるいは「意思決定」を行うための方法として，モデル比較やNHSTといった方法がある。
	\item[統計環境Rによる実践] 統計環境Rを利用して，理論的な理解だけでなく実際に計算ができることも身につけるべき技術である。
\end{description}

\subsection*{標準シラバスとの対応}
各コマの「公認心理師標準シラバスにおける位置づけ」と「心理統計法標準カリキュラム(日本心理学会)における位置づけ」には，それぞれ次の標準シラバスとの対応を示した。
\begin{description}
	\item[公認心理師標準シラバス] 公益社団法人日本心理学会・公認心理師養成大学教員連絡協議会「公認心理師大学カリキュラム 標準シラバス」(2018年8月22日版)。科目番号，大項目，中項目の順に示す。
	\item[心理統計法標準カリキュラム] 日本心理学会 心理統計法標準カリキュラム作成小委員会「ベシスタ コマシラバス案」(2022年8月5日更新)。回と細目の番号で示す。\url{https://sites.google.com/view/jpa-psychometrics/home}
\end{description}

\newpage
\tableofcontents
\newpage

\section{心理学と統計法}
\label{b01_Introduction}
\subsection{授業内容}
\subsubsection{概要}
心理学を学ぶに当たって，なぜ統計学を学ばなければならないのかを考える。そのためには，心理学がサイエンスであること，サイエンスであるとは客観的であること，積み重ねが必要であることを知る必要があり，そのために数量的表現が有用であることを理解する。
\subsubsection{コマ主題細目}
\begin{description}

	\item[心理学とはどういう学問か] 心理学は初等教育の中には取り入れられていないにもかかわらず，心や気持ちといった用語は日常的に頻繁に用いられ，かつ重要視されているものでもある。「心理学の過去は長いが歴史は短い」といわれるが，どのようにして近代科学の仲間入りをしたのかを理解するため，サイエンスとしての枠組みを改めて問い直す。
	      \begin{flushright}
		      $\to$\textcite{Miura201705}のP.1--7.
	      \end{flushright}

	\item[心理学の2つの流派] 心理学には基礎と応用，という区別がされることがあるが，内容やその方法論的な区分から行けば実証主義的アプローチと了解主義的アプローチの2つに分類する方が良い。また基礎と応用という区分では，基礎が先に，下にあって応用が後に，上にくるような印象を持つが，それらは理学と工学の違いのように，併存して然るべきものである。
	      \begin{flushright}
		      $\to$\textcite{Shimoyama2001}による研究方法の分類
	      \end{flushright}

	\item[近代科学の特徴] 心理学は科学である，というならばその科学とはそもそもなんだろうか。ここでは近代科学の特徴として実験，数学的現象主義，機械論，要素還元主義をあげる。ニュートンやガリレオが行なったのは，関数による数学的現象主義であり，心理学そのものは基本的にその枠組みの中に位置する。すなわちアウトカム変数に対して心というモデルを立てて，そのモデルの正しさを検証することが1つの目的である。アウトカム変数は数量化されていなくてもよいが，数量化されていた方が正確性，客観性，比較可能性，要約可能性に優れており，かつ統計モデルの活用により表面的に現れる数字から見えない情報を引き出すことが可能になっている。
	      \begin{flushright}
		      $\to$\textcite{Miura201705}のP.79--81.
		      $\to$数値化の意味については\textcite{Kawabata201408}のP.1-6
	      \end{flushright}

	      % \section{心理学のいとなみ} 心理学はそもそもがモデル論であるが，これをとくに数値データと対応させ，実際に検証しながら研究は進められる。心理学のリサーチクエスチョンの立て方は，効果の有無を考える場合(平均的因果効果の推論)とメカニズムを考える場合(モデリングアプローチ)の二種類が存在するが，いずれも標本の特徴から母集団の特徴を推論するという枠組みでは同じである。これらを踏まえて，心理統計的アプローチの必要性を理解する。

\end{description}
\subsubsection{キーワード}
\begin{itemize}
	\item 心理学における2つの検証方法
	\item 近代科学としての心理学
	\item 数学的現象主義
\end{itemize}
\subsection{授業情報}
\paragraph{コマの展開方法}
講義
\paragraph{公認心理師標準シラバスにおける位置づけ}
科目番号4 心理学研究法；①心理学における実証的研究法(量的研究及び質的研究)；A 科学と実証 に相当
\paragraph{心理統計法標準カリキュラム(日本心理学会)における位置づけ}
第1回 オリエンテーション，心理学と統計，変数とデータ(位置づけ・到達目標) に相当
\subsubsection{予習・復習課題}
\paragraph{予習}
心理学入門の入門書として，\textcite{Daihiko201507,Michimata200903}などを一読しておくことで，これから4年間かけて学ぶ心理学についての事前のイメージを再構築すると良い。ほかにも，日本心理学会が一般向けに出版している「心理学ワールド」が読みやすい入門雑誌であろう。あるいは心理学についてのイメージ調査については\textcite{Kikuchi2018}が詳しく，心理学の歴史についての専門書としては\textcite{Ohashi201610}なども参考にすると良い。
\paragraph{復習}
本学における心理学の学習が始まる前に持っていた，心理学に対する事前のイメージをいったんわすれ，まずは広く心理学の分野について眺めた上で，入学を志望したときに持っていた心理学についてのイメージを改めて位置づけてもらいたい。心理学のより細かい専門領域に対する位置付けを把握した上で，心理学統計法のような研究方法が統一され共有されていること，とくに初年時の必須科目として位置付けられていることの意義を自分なりに見出せるように，ノートなどに授業の感想と現段階での今後の展望を書き留めるなどすると良い。

\section{心理学と計算機}
\label{s01_Introduction}
\subsection{授業内容}
\subsubsection{概要}
心理統計の授業では，多くの数値を一気に扱うため計算機を用いる。Excel/Numbers/Calcsのような表計算ソフトでは不十分なため，Rと呼ばれる統計専門のソフトウェアを用いる。Rの使い方については次回以降に詳細が論じられるが，それに先立って計算機の基礎的な概念，操作方法を理解する。
\subsubsection{コマ主題細目}
\begin{description}

	\item[電子計算機の基礎] 計算機には5つの機能がある。すなわち，入力，出力，演算，制御，記憶装置とよばれるもので，対応するデバイスの組み合わせで計算機が構成される。また，BIOS，OS，アプリケーションの違いなど，計算機の内部構造についての理解は，人間をコンピュータとみなして考えるとき，メタファとして有用である。
	
	\item[情報の単位] 電子計算機はすべて0/1の情報として扱う。さまざまな単位と，対応する装置ではどの程度の大きさを扱えるのかについて，基本的な知識を身につける。ビット，バイト，キロバイトからテラバイトまでの関係性や，テキストデータと画像・動画データの容量の違いなどを理解することで，扱うデータの規模感を把握する。

	\item[ファイルの種類と拡張子] コンピュータで扱うファイルには様々な種類があり，拡張子によって区別される。特に統計解析で扱うことの多いテキストファイル(.txt)，CSV形式(.csv)，Excelファイル(.xlsx)の違いを理解し，それぞれの特徴と用途を把握する。また，Rで扱うデータファイル(.RData)やスクリプトファイル(.R)についても概説する。

	\item[ファイルの場所] コンピュータ内でのファイル管理の仕組みとして，ディレクトリ(フォルダ)階層構造について学ぶ。カレントディレクトリ，相対パス，絶対パスなどの概念を理解し，Rを使った解析においてファイルの読み込みや保存をする際に必要となる基礎知識を習得する。クラウドストレージの利用についても触れ，バックアップの重要性を理解する。

	\item[キーボードとショートカット] 効率的なコンピュータ操作のために，キーボードの配置やタイピングの基本，よく使うショートカットキー(コピー・ペースト，検索，保存など)について学ぶ。特にRやRStudioでの作業に役立つキーボード操作を中心に解説し，実際に操作しながら習得していく。


\end{description}
\subsubsection{キーワード}
\begin{itemize}
	\item コンピュータの基本構造(入力・出力・制御・演算・記憶)
	\item ファイル形式とデータ管理
	\item 効率的なPC操作
\end{itemize}
\subsection{授業情報}
\paragraph{コマの展開方法}
講義
\paragraph{公認心理師標準シラバスにおける位置づけ}
該当なし
\paragraph{心理統計法標準カリキュラム(日本心理学会)における位置づけ}
該当なし
\subsubsection{予習・復習課題}
\paragraph{予習}
大学のPCルームやメディアセンターで利用できるコンピュータの基本操作方法を確認しておくこと。自分のノートPCを授業に持参する場合は，事前にOSのバージョン，ストレージの空き容量，メモリ容量を確認しておくこと。また，普段使用しているデバイス(スマートフォンやタブレット)と比較して，PCではどのような操作が異なるかについて考えてみると良い。
\paragraph{復習}
授業で紹介したファイル操作(作成・保存・移動・削除)を自分のPCで実践してみること。特に，新しいフォルダを作成し，その中にテキストファイルを作成・保存する練習をすること。また，キーボードショートカット(Ctrl+C, Ctrl+V, Ctrl+Z, Ctrl+S など)を意識的に使用する習慣をつけるよう心がけること。次回の授業までに，Rとそのインターフェースとなるソフトウェア(RStudio)をインストールしておくこと。

\section{統計環境Rの導入}
\label{b03_introR}
\subsection{授業内容}
\subsubsection{概要}
携帯電話やタブレットが誰でも身近に利用できるようになってきてはいるが，それらはあくまでも受動的にサービスを利用する媒体であり，執筆や計算などユーザが能動的に電子機器を利用するにあたっては，パソコンの利用は不可欠である。
本講義だけでなく，心理学の研究をする上でも当然パソコンの利用は必須であり，とくに統計的データ処理をするにあたってはより深く知る必要がある。また，ExcelやNumbersなどの表計算でも記述統計量やグラフの作図はできるが，それ以上の「複数の変数を扱った多変量解析」や「群間の平均値の差を細かく検証する」といった心理統計的な作業には不十分なのである。
そこで，より専門的なツールとして統計環境Rを導入する。統計環境Rは専修大学人間科学部心理学科が公式に採用している統計環境であり，これを活用する総合開発環境であるRStudioの基本的な使い方，結果の再現性を担保するための工夫などについても言及する。それに先立って，コンピュータの基礎知識についても解説する。
\subsubsection{コマ主題細目}
\begin{description}

    \item[統計環境R/RStudioの準備] 心理学科では基本的に統計環境Rを共通のツールとして用いる。統計環境Rはフリーソフトウェアであり，誰でも無償で利用できる環境である。また，R単体では簡素な言語的やりとりができるに過ぎない。ファイルの操作や描画機能，記録，参照，管理など統計言語を扱うに当たって必要な総合的環境としてRStudioを併用することで，その利便性は飛躍的に向上する。まずはR/RStudioを導入し，基本的な画面構成やプロジェクトによる管理ができることを確認する。
	      \begin{flushright}
		      $\to$ Rについての入門書は多くあるがここでは\textcite{kosugi2019}のP.1--7を挙げておく。同じくRStudioを使った入門書は多くあるがここでは\textcite{kosugi2019}のP.7--24を挙げておく。
	      \end{flushright}

	\item[Rの基本操作]環境が整ったところで，Rの基本的な操作を見ていく。まずはRスクリプトをつかって，四則演算や関数計算，データの形式やパッケージの導入と利用について理解する。スクリプトペインとコンソールペインの違いを理解し，またスクリプトとして清書しておくことで計算手続きの記録がつけられることの利点を理解する。

	\item[基本的なデータ型と演算] RではNumeric(数値)，Character(文字列)，Logical(論理値)などの基本的なデータ型があり，それぞれの特徴を理解する。またベクトル，行列，データフレームといったデータ構造の違いを学び，基本的な統計量の算出方法を実践する。
          \begin{flushright}
              $\to$\textcite{kosugi2019}のP.25--40
          \end{flushright}

\end{description}
\subsubsection{キーワード}
\begin{itemize}
    \item R/RStudio
    \item スクリプトとコンソール
    \item データ型とデータ構造
    \item パッケージ管理
\end{itemize}
\subsection{授業情報}
\paragraph{コマの展開方法}
講義
\paragraph{公認心理師標準シラバスにおける位置づけ}
科目番号5 心理学統計法；②統計に関する基礎的な知識；C エクセル，R，SPSS入門(1) に相当
\paragraph{心理統計法標準カリキュラム(日本心理学会)における位置づけ}
該当なし
\subsubsection{予習・復習課題}
\paragraph{予習}
事前にR(https://www.r-project.org/)とRStudio(https://posit.co/download/rstudio-desktop/)をダウンロードしてインストールしておくこと。インストールに問題がある場合は，授業開始前に教員やTAに相談すること。

\paragraph{復習}
授業で学んだR/RStudioの基本操作(四則演算，関数の使用，データの入力など)を実際に自分のPCで繰り返し練習すること。特に，スクリプトを用いた操作に慣れること。また，簡単な数値データを用いて平均値や標準偏差を算出する練習を行うこと。
\section{数値データの種類}
\label{b02_Number_and_data}

\subsection{授業内容}
\subsubsection{概要}
目に見えない心という対象を研究するために，言語や行動によって得られた反応を記録しなければならず，記録されたものは集計，分析のために数値化される。数字やその統計的分析結果が，対象を必ずしも反映していないという直観に反するかもしれないが，数字の限界は言語の限界と同じであり，言い換えれば論述・論考・論証の限界に等しいすべての科学的領域をカバーする。ここでは数値化することで何が得られるか，その長所を理解するとともに，データの特性に応じた適切な分析方法を選択するための基礎知識を身につける。
\subsubsection{コマ主題細目}
\begin{description}
	\item[データと変数] 個人の特徴を記述するにあたって，心理統計ではすべての情報を数値化する。これらの数字は人によって変わりうるので，一般に変数と呼ぶ。また，目に見える変数を観測変数やアウトカムとよぶが，目に見えない心の状態をデータから取り出したりする場合は，それを潜在変数と呼んで区別する。変数の性質を理解することは，適切な統計的手法を選択する上で重要である。
	      \begin{flushright}
		      $\to$ データの分類については，\textcite{Yamada2004}のP.18--21。
	      \end{flushright}

	\item[計算機と統計の観点による分類] 統計ソフトウェアでは，データの特性に応じて異なる型を使用する。バイナリ(二値)データは0/1の二つの値のみを取る変数であり，性別や有無といった情報がこれに該当する。カテゴリカルデータは複数のカテゴリーに分類される変数で，数値的な順序を持たない。オーディナルデータは順序を持つカテゴリカルデータであり，アンケートの「全く同意しない」から「強く同意する」などの尺度がこれに相当する。連続変数は身長や体重などの無限の値をとりうるデータであり，コンピュータ上では離散的に扱われるが理論上は連続的である。それぞれのデータ型はRなどの統計ソフトウェアで異なる方法で処理され，適切なデータ型の指定が分析の前提となる。
	      \begin{flushright}
		      $\to$ Rにおけるデータ型については\textcite{Baba2020}のP.65--74を参照。
	      \end{flushright}

	\item[尺度水準による分類] すべての数字が四則演算が可能かと言われると，そうではない。尺度水準とは，数値が持つ情報や性質の段階を表す言葉であり，Steavens,S.が提唱した4段階の区別がとくに重要である。名義尺度(nominal scale)は単なる分類を表し，順序尺度(ordinal scale)は順序関係を表す。間隔尺度(interval scale)は等間隔性を持ち，比率尺度(ratio scale)は絶対的なゼロ点を持つという特徴がある。これらの尺度水準の違いは，適用可能な統計的手法を決定するため，得られた数字データがどの水準に該当するかを判別することができるようにならなければ，誤った統計処理をすることになるため注意が必要である。
	      \begin{flushright}
		      $\to$\textcite{Kawabata201408}のP.9--16，あるいは\textcite{Yamada2004}のP.22-25.
	      \end{flushright}

	\item[コンピュータデータ型と尺度水準の関係] コンピュータにおけるデータ型と心理統計学における尺度水準は緊密に関連している。例えば，名義尺度はカテゴリカルデータとして，順序尺度はオーディナルデータとして扱われる。間隔尺度と比率尺度は通常，連続データとして扱われるが，その解釈と扱いには重要な違いがある。Rなどの統計ソフトウェアでは，これらの違いを適切に指定することで，データの特性に合った分析が可能となる。実際の研究においては，測定方法や理論的背景を考慮して，データの性質を適切に判断することが求められる。

\end{description}

\subsubsection{キーワード}
\begin{itemize}
	\item 観測変数と潜在変数
	\item バイナリ・カテゴリカル・オーディナル・連続データ
	\item 確率分布とデータ型
	\item 尺度の四水準(名義・順序・間隔・比率)
\end{itemize}

\subsection{授業情報}
\paragraph{コマの展開方法}
講義と実習
\paragraph{公認心理師標準シラバスにおける位置づけ}
科目番号4 心理学研究法；①心理学における実証的研究法(量的研究及び質的研究)；C 実証の手続き に相当
\paragraph{心理統計法標準カリキュラム(日本心理学会)における位置づけ}
第1回 オリエンテーション，心理学と統計，変数とデータ 細目1・2 に相当
\subsubsection{予習・復習課題}
\paragraph{予習}
一般に，データに基づいて(Ebidence Based)論証・考察することが主観的な思い込みから脱却する第一歩である。尺度水準については心理学や心理統計の入門的テキストでは必ず触れられる項目であり，\textcite{Miura201705,Kawabata201408,Yamada2004,kosugi2018,kosugi2019}などを手に取って確認しておくと良い。また，Rを用いたデータ操作の基本について，前回の復習を行っておくこと。
\paragraph{復習}
身の回りにあるさまざまな数字に注意を払い，それらがどの尺度水準の情報を持っているか，またコンピュータ上ではどのようなデータ型で扱うべきかを考えることで理解を深める。例えば，学生証番号，授業の成績評価(S/A/B/C/D)，テストの点数，体重などを考え，それぞれがどの尺度水準に該当するか，Rではどのデータ型で表現すべきかを整理してみること。また，各種の心理尺度(例：リッカート尺度)がどの尺度水準に該当するかについても考察し，適切な統計的処理について考えてみると良い。次回の授業に向けて，Rで実際にさまざまなデータ型を扱う練習をしておくこと。
\section{記述統計量とRによる演習}
\label{b05_descriptives_R}

\subsection{授業内容}
\subsubsection{概要}
統計の基本は集めたデータを集計し，要約統計量で表現することである。手元にあるデータの詳細を記述するため，これは記述統計(descriptive statistics)と呼ばれる。いかなる進んだ分析を行うに当たっても，必ず可視化および記述統計による特徴の把握を欠かすことはできない。また要約統計量はその性質上，必ず有意味な情報を捨て去って代表的な性質を表すものであり，単一の指標だけで考えることは危険であることを理解し，活用できるようにする。本コマでは記述統計の概念を学ぶとともに，Rを用いた実践的な演習を通じてデータの要約と可視化の技術を身につける。
\subsubsection{コマ主題細目}
\begin{description}
	\item[中心化傾向の指標] 散らばりを持つデータの中心についての要約統計量として，平均値(mean)，中央値(median) ,最頻値(mode)の3つがあげられる。それぞれの特徴および算出方法，その際の注意点についても理解する。また，計算や表現の利便性のため，$\sum$という総和記号を用いる。平均値は外れ値の影響を受けやすく，中央値はその影響を受けにくいという特性の違いから，データの性質に応じた適切な指標の選択が重要であることを理解する。
	      \begin{flushright}
		      $\to$ \textcite{Yamauchi201001}のP.27--35,\textcite{Yamada2004}のP.30--33，\textcite{Kawabata201408}のP.23--26.
	      \end{flushright}

	\item[散らばりの指標] データにとって重要なのは中心化傾向の指標だけではなく，むしろその散らばりについての指標である分散(variance),標準偏差(standard deviation)こそ重要である。加えて最大値(max)，最小値(min) ,四分位偏差(quantile)，パーセンタイル(percentile)，範囲(range)などが散らばりを表す指標として用いられる。ここではそれぞれの特徴および算出法について理解する。特に標準偏差は平均値からの距離の二乗平均の平方根として定義され，データの散らばり具合を把握する上で最も一般的に用いられる指標であることを理解する。
	      \begin{flushright}
		      $\to$ \textcite{Yamauchi201001}のP.37--48,\textcite{Yamada2004}のP.34--37，\textcite{Kawabata201408}のP.26-33.
	      \end{flushright}

	\item[Rによる集計] 実際の研究では大量のデータを扱うため，手計算ではなくコンピュータを用いた集計が不可欠である。Rでは外部データファイルの読み込みからデータの概要把握までを効率的に行うことができる。ここではCSVファイルの読み込み方法を確認し，読み込んだデータの構造を確認するための\texttt{str()}関数や，基本的な要約統計量を一度に出力する\texttt{summary()}関数の使用法を学ぶ。データの読み込みにおけるエンコード指定や，ヘッダの取り扱いなど，実践的なデータ処理の基礎を身につける。
	      \begin{flushright}
		      $\to$ \textcite{kosugi2019}のP.60--70，\textcite{Kinosady2021}の第2章。
	      \end{flushright}

	\item[Rによる統計量の算出] Rには記述統計量を算出するための様々な関数が用意されている。代表的なものとして\texttt{mean()}, \texttt{median()}, \texttt{sd()}, \texttt{var()}, \texttt{min()}, \texttt{max()}, \texttt{quantile()}などがあり，これらを用いてデータの特徴を数値的に把握する方法を学ぶ。欠測値(\texttt{NA})の取り扱いについても，実データを用いて学習する。
	      \begin{flushright}
		      $\to$ \textcite{kosugi2019}のP.73--85，\textcite{Baba2020}のP.156--182。
	      \end{flushright}

	\item[Rによるデータの可視化] データの分布を視覚的に理解するために，Rの標準プロット関数を用いたヒストグラム(\texttt{hist()})と箱ひげ図(\texttt{boxplot()})の描画方法を学ぶ。特にヒストグラムでは階級幅の設定による表示の違いや，箱ひげ図における外れ値の表示方法についても理解する。これらの可視化を通じて，データの分布形状や歪み，外れ値の有無などを把握し，適切な記述統計量の選択につなげる方法を身につける。また，複数グループの比較を行うための工夫についても学ぶ。
	      \begin{flushright}
		      $\to$ \textcite{Kinosady2021}の第4章，\textcite{kosugi2019}のP.90--105。
	      \end{flushright}

\end{description}
\subsubsection{キーワード}
\begin{itemize}
	\item 平均値，中央値，最頻値
	\item 分散，標準偏差，四分位偏差，範囲
	\item データの読み込みと前処理
	\item R関数による記述統計量の算出
	\item ヒストグラム，箱ひげ図による可視化
\end{itemize}
\subsection{授業情報}
\paragraph{コマの展開方法}
講義と演習：実際のPCを使う実習を伴った授業を行う
\paragraph{公認心理師標準シラバスにおける位置づけ}
科目番号5 心理学統計法；①心理学で用いられる統計手法；A 記述統計：代表値と散布度／科目番号5 心理学統計法；②統計に関する基礎的な知識；E エクセル，R，SPSS入門(3) に相当
\paragraph{心理統計法標準カリキュラム(日本心理学会)における位置づけ}
第2回 1変数の標本統計量 細目1・2・3，第1回 オリエンテーション，心理学と統計，変数とデータ 細目3 に相当
\subsubsection{予習・復習課題}
\paragraph{予習}
統計分野は2012年(平成24年)より，高校数学の必修である数学Iで新たに扱われることとなっているため，受講生諸君はすでに算術平均や分散について一度は習っていることと思われる。基本的な統計量の計算方法について自信がない場合は，高校数学のテキストを読み返しておくこと。また，前回までにインストールしたRとRStudioの基本操作を復習しておくこと。
\paragraph{復習}
授業で扱ったサンプルデータを用いて，Rmarkdownで記述統計レポートを作成すること。具体的には，(1)データの読み込み，(2)基本的な記述統計量の算出，(3)ヒストグラムと箱ひげ図の作成，(4)結果の解釈という流れで，Rのコードと結果，そして結果の解釈を含むレポートをRmarkdownで作成し提出すること。なお，このレポートの提出は前期の単位認定にあたっての必要条件であり，基準に満たないレポートは再提出を求める。前期試験の前にパスすることが，前期末試験を受験するための前提条件であることを承知しておくこと。

また，記述統計の基本概念については，\textcite{Yamada2004}や\textcite{Kawabata201408}などのテキストを参照し，理解を深めておくこと。特に，データの性質に応じた適切な記述統計量の選択について考察しておくとよい。Rの操作については\textcite{kosugi2019}が初学者向けにわかりやすく解説しているので，演習の復習に活用するとよい。
\section{相関係数とRによる実践}
\label{b06_Corr}

\subsection{授業内容}

\subsubsection{概要}
これまでは1変数の記述統計について学んできたが，本コマからは複数の変数間の関係性について考察する。変数同士の関係性を数値化して表す指標として，共分散や相関係数について学ぶ。また，変数の尺度水準に応じて適切な関連指標が異なることを理解し，それぞれの特徴と適用条件について理解する。相関係数は直線的関係についての指数であることに注意し，散布図による可視化の重要性を理解する。実際のデータを用いてRによる計算と可視化の演習を行うことで，相関分析の実践的技能を身につける。

\subsubsection{コマ主題細目}
\begin{description}
	\item[二変数間関係の指標] 変数同士の関係を記述する方法として，クロス集計表，共分散，相関係数などがある。クロス集計表はカテゴリカルデータの関連を表現する基本的な方法である。共分散は二つの変数がどのように連動して変化するかを数値化したものだが，変数の単位に依存するという欠点がある。そこで単位の影響を取り除くために標準化したものが相関係数であり，-1から+1の範囲で二変数間の線形関係の強さを表す。これらの基本的な考え方と計算方法について理解する。
	      \begin{flushright}
		      $\to$ \textcite{Yamauchi201001}のP.62--68,\textcite{Yamada2004}のP.44--50，\textcite{Kawabata201408}のP.48--52.
	      \end{flushright}

	\item[尺度水準に応じた相関係数] 変数の尺度水準によって適切な関連指標は異なる。名義尺度の関連はカイ2乗値やファイ係数，クラメールのV係数などで表される。二値データ間の関連を表すテトラコリック相関，順序尺度間の関連を表すポリコリック相関，順序尺度と連続変数の関連を表すポリシリアル相関などがある。連続変数間の関連を表す指標としては，ピアソンの積率相関係数が最も一般的であるが，線形関係以外も検出できるスピアマンの順位相関係数やケンドールのタウ，さらに最近では非線形関係も捉えられるMIC(Maximum Information Coefficient)なども開発されている。これらの指標の特徴と適用条件を理解し，データの性質に応じた適切な指標を選択できるようになる。
	      \begin{flushright}
		      $\to$ 様々な相関係数について\textcite{Kawabata201408}のP.53--58，\textcite{Yamada2004}のP.51--59，MICやテトラコリック相関などカテゴリカルな相関係数については\textcite{TDE2022}の2章を参照。
	      \end{flushright}

	\item[Rによる相関分析の実践] 相関係数の計算と検定はRを用いて簡単に行うことができる。\texttt{cor()}関数を用いた相関係数の計算方法について学ぶ。また，\texttt{polycor}パッケージを用いたポリコリック相関やポリシリアル相関の計算方法，\texttt{minerva}パッケージを用いたMICの計算方法についても触れる。実データを用いた演習を通じて，相関分析の実践的な技術を身につける。
	      \begin{flushright}
		      $\to$ Rによる相関分析については\textcite{kosugi2019}のP.110--120を参照。
	      \end{flushright}

	\item[相関の落とし穴と可視化の重要性] 相関係数だけを見て判断することの危険性について理解する。同じ相関係数の値でも，散布図のパターンは大きく異なる可能性がある。この点を示すための有名な例として，アンスコムの例や，より現代的な\texttt{datasauRus}パッケージに含まれるDinosaurus datasetがある。これらの例を通じて，相関分析を行う際には必ず散布図を描いてデータの分布を確認することの重要性を学ぶ。また，外れ値の影響や非線形関係，層別の問題など，相関分析における様々な落とし穴について理解し，それらを回避するための方法について学ぶ。
	      \begin{flushright}
		      $\to$ \verb|datasauRus|パッケージには，記述統計が同じだがプロットがまったく異なる13のケースが含まれていて参考になる。可視化の重要性については\textcite{Kieran2018}を参照。
	      \end{flushright}
\end{description}

\subsubsection{キーワード}
\begin{itemize}
	\item 共分散と相関係数
	\item カイ2乗値，ファイ係数，テトラコリック相関，ポリコリック相関
	\item ピアソンの積率相関係数，スピアマンの順位相関，MIC
	\item 散布図による可視化
\end{itemize}

\subsection{授業情報}
\paragraph{コマの展開方法}
講義と演習
\paragraph{公認心理師標準シラバスにおける位置づけ}
科目番号5 心理学統計法；①心理学で用いられる統計手法；B 記述統計：相関 に相当
\paragraph{心理統計法標準カリキュラム(日本心理学会)における位置づけ}
第3回 2変数の関係 細目3・4・5 に相当
\subsubsection{予習・復習課題}
\paragraph{予習}
分散，標準偏差，標準化の手続きについて再確認しておくこと。分散は自分自身との共分散という形で理解することができるし，共分散を標準化したものが相関係数になるからである。また，Rの基本操作とデータ読み込みの方法について復習しておくこと。
\paragraph{復習}
授業で扱ったサンプルデータを用いて，以下の課題を行い，Rmarkdownファイルにまとめて提出すること。
1. 複数の変数間の相関係数を計算し，相関行列を作成する
2. 相関係数のヒートマップを作成する
3. 散布図を描き，相関係数の値と視覚的な関係性を比較考察する
4. \texttt{datasauRus}パッケージの例を再現し，同じ相関係数でも異なるパターンが生じることを確認する
5. 身近な例から二つの変数を選び，それらの関係を相関分析によって検討する

また，相関関係と因果関係の違いについて理解を深めるため，日常生活や報道などで見られる「相関関係を因果関係と誤解している例」を一つ挙げ，なぜそれが誤りであるかを説明すること。これはデータリテラシーの問題でもあり，統計的思考の基本でもある。
\section{回帰分析の基礎}
\label{b07_Regression}

\subsection{授業内容}

\subsubsection{概要}
前回学んだ相関係数は変数間の関連性を示すものであるが，一方の変数から他方の変数を予測・説明する枠組みとして回帰分析がある。本コマでは最も基本的な単回帰分析について学び，予測モデルの構築と評価の方法を理解する。また，相関関係と因果関係の違いについて考察し，因果推論の基礎的な考え方を学ぶ。古典的テスト理論から測定に含まれる誤差を知り，誤差が正規分布に従っていると考えることで，平均の位置を推定することができるようになった考え方を拡張し，他の変数との関係から平均の位置が関数で表現される最も単純な場合である線形モデルについて考える。

\subsubsection{コマ主題細目}
\begin{description}
	\item[予測と線形モデル] 変数間関係を記述する最も単純な形として，一次関数を用いたものを考える。現象に対して関数で記述する数学的現象主義，変数を説明変数と被説明変数として捉えること，予測に誤差が加わって現象となることなど，これまでの議論をデータと数式に置き換えただけであることをしっかりと理解する。さらに単回帰分析の数式的記述について理解する。とくにデータ($X_i,Y_i$)，予測値($\hat{Y}$)，誤差($e_i$)，回帰係数($b_0,b_1$)として表される記号の細部にまで注意を払い，その意味を理解する。
	      \begin{flushright}
		      $\to$ \textcite{Toyoda201706}のP.27--28.
	      \end{flushright}

	\item[最小二乗法] 最小二乗法によって推定値が算出できることを理解する。数式の細かい展開はここでは行わず，最小化する目的関数の直観的理解に止める。また算出された係数(傾きと切片)がどのような意味を持つのか改めて理解する。傾きは説明変数が1単位変化したときの被説明変数の変化量を表し，切片は説明変数がゼロのときの被説明変数の値を表す。これらの係数の解釈と実践的な意味について考察する。
	      \begin{flushright}
		      $\to$　\textcite{Toyoda201706}のP.30
	      \end{flushright}

	\item[残差と決定係数] 回帰係数が計算されれば，実際のデータに対して予測値を計算することができ，予測がどの程度適合していたかを検証することができる。予測値と被説明変数との差分を残差(residuals)と呼ぶが，この両者の相関係数を重相関係数，その二乗したものを決定係数($R^2$)と呼んでモデルの適合度を測る指標にすることを理解する。決定係数は被説明変数の分散のうち，モデルによって説明される割合を表し，モデルの説明力を示す重要な指標である。
	      \begin{flushright}
		      $\to$\textcite{Nojima201903}のP.58-60.
	      \end{flushright}

	\item[相関関係と因果関係] 相関関係があるところに，因果的な説明をしてしまいがちであるが，因果的な説明をするためには相関以上に注意が必要である。因果関係を見るために必要な，時間的先行性，他の原因がないこと，他の結果に至らないことなどの諸条件を確認する。回帰分析は変数間の関係を数理モデルとして表現する強力なツールであるが，それだけでは因果関係を示すには不十分であることを理解する。回帰係数が因果効果として解釈できる条件と，そうでない場合の違いについて理解する。
	      \begin{flushright}
		      $\to$ 相関と因果の違いについては，\textcite{Toyoda1992}に簡潔にまとめられている。
	      \end{flushright}

\end{description}
\subsubsection{キーワード}
\begin{itemize}
	\item 一次関数と線形モデル
	\item 予測値，誤差，回帰係数
	\item 最小二乗法と決定係数
	\item 相関関係と因果関係
\end{itemize}
\subsection{授業情報}
\paragraph{コマの展開方法}
講義と演習
\paragraph{公認心理師標準シラバスにおける位置づけ}
科目番号5 心理学統計法；①心理学で用いられる統計手法；C 記述統計：回帰／科目番号4 心理学研究法；②データを用いた実証的な思考方法；A 相関関係から因果関係へ に相当
\paragraph{心理統計法標準カリキュラム(日本心理学会)における位置づけ}
第4回 相関と回帰 細目1・2・3・4，第5回 相関と因果 細目1 に相当
\subsubsection{予習・復習課題}
\paragraph{予習}
前回学習した相関係数の概念と計算方法について復習しておくこと。また，一次関数の基本的な性質(傾きと切片の意味)についても確認しておくと良い。
\paragraph{復習}
\textcite{Yamauchi201001}のP.62--81は相関係数と回帰分析を1つの章にまとめて解説しているので参考にすると良い。最小二乗法による回帰係数の算出については，偏微分連立方程式を用いれば簡単に計算できるが，偏微分連立方程式とは何かに遡って解説している本として\textcite{kosugi2018}のP.104--112がある。

また，授業で学んだ内容を応用して以下の課題に取り組むこと：
\begin{enumerate}
	\item 実際のデータセットを用いて単回帰分析を実施し，回帰係数と決定係数を算出する
	\item 得られた回帰式を用いて予測を行い，残差をプロットする
	\item 身近な例から「相関関係と因果関係の混同」の事例を見つけ，なぜそれが因果関係とは言えないのかを説明する
	\item 因果関係を主張するためには，相関関係に加えてどのような条件や証拠が必要かを考察する
\end{enumerate}

\section{重回帰分析と変数間の影響}
\label{b08_MultipleRegression}

\subsection{授業内容}
\subsubsection{概要}
前回は単回帰分析について学んだ。今回は複数の説明変数を用いる重回帰分析へと発展させる。回帰分析によって，データにモデルを当てはめて係数の推定値を得るということを行ったが，そのこととモデルが適していたかどうかは別である。この点を理解するためには，モデルの適合度について考える必要がある。モデルが適合しているということを踏まえた上で，説明する変数が増える重回帰分析へと議論は展開する。ここでは説明変数の読み取り方への注意を促すべく，偏相関係数の考え方を解説し，偏回帰係数，標準化偏回帰係数などの用語について理解する。

\subsubsection{コマ主題細目}
\begin{description}
	\item[重回帰分析の基本] 説明変数が複数に増えた場合の回帰分析は，重回帰分析(Multiple Regression Analysis)と呼ぶ。数式的な表現としては項が増えるだけであるが，データのプロットを考えると三次元以上の回帰平面に拡張されていること，それでも面は湾曲していない，すなわち線形関係であることを理解する。単回帰分析と同様に，最小二乗法によって回帰係数を推定するが，行列計算を用いることで効率的に解が求められることにも触れる。
	      \begin{flushright}
		      $\to$ \textcite{kosugi2019}のP.168,図10.7の回帰平面図を参照
	      \end{flushright}


	\item[部分相関と偏相関と偏回帰係数] 重回帰分析における係数の解釈のために，部分相関と偏相関の概念を理解する。回帰分析によって，被説明変数の分散が説明される部分と残差の部分に分離させられること，残差と説明変数は相関しないことを確認し，残差は説明変数の影響を除外したものと考えることができる。これは条件を必ずしも統制できない事象において，統計モデル的に統制していることでもある。偏相関係数は他の変数の影響を取り除いた2変数間の純粋な関連の強さを表す指標であり，変数間の直接的な関係を評価する上で重要な概念である。
	      \begin{flushright}
		      $\to$ \textcite{Kawabata201408}のP.59-65，\textcite{Yamada2004}のP.60-65。
	      \end{flushright}

	\item[偏回帰係数と解釈] 部分相関の説明において，残差が分離されることを見てきたが，残差同士の相関関係はとくに偏相関ということができる。これは共通する変数からの影響を除いた変数間関係であり，これを見るだけでも変数間関係を見ることにつながる。とくに残差変数を使っての回帰分析を行うとき，得られる係数が偏回帰係数である。重回帰分析における回帰係数はこのように条件付きで考えなければならない。偏回帰係数は「他の説明変数の値を一定に保った場合に，その説明変数が1単位変化したときの被説明変数の変化量」として解釈できることを理解する。また変数を増やすことで重相関係数が増加することも理解する。また，複数の説明変数間に強い相関がある場合，多重共線性(multicollinearity)の問題が生じる可能性があることを理解する。多重共線性があると，回帰係数の推定値が不安定になり，解釈が困難になるという問題が発生する。VIF(Variance Inflation Factor)などの診断指標や，多重共線性が検出された場合の対処法についても簡単に触れる。
	      \begin{flushright}
		      $\to$\textcite{kosugi2018}のP.60--63.多重共線性については\textcite{Shimizu201705}のP.86--89を参照。
	      \end{flushright}

	\item[標準化された係数] 回帰係数については，素点による回帰係数が理解しやすいが，単位に左右されるものでもあるので相対比較をする場合はすべての変数を標準化した標準化解を用いた方がわかりやすい。標準化解を用いても適合度は変化しないことなど留意点とともに有用性を理解する。
	      \begin{flushright}
		      $\to$\textcite{kosugi2018}のP.66--67，
	      \end{flushright}

\end{description}
\subsubsection{キーワード}
\begin{itemize}
	\item 重回帰分析と回帰平面
	\item 交絡要因と統計的制御
	\item 部分相関と偏相関
	\item 偏回帰係数と標準化偏回帰係数
	\item 多重共線性
\end{itemize}
\subsection{授業情報}
\paragraph{コマの展開方法}
講義と演習
\paragraph{公認心理師標準シラバスにおける位置づけ}
科目番号5 心理学統計法；①心理学で用いられる統計手法；C 記述統計：回帰／科目番号4 心理学研究法；②データを用いた実証的な思考方法；D 複雑な心理事象のモデリング に相当
\paragraph{心理統計法標準カリキュラム(日本心理学会)における位置づけ}
第13回 重回帰分析 細目1・2・4，第5回 相関と因果 細目2 に相当
\subsubsection{予習・復習課題}
\paragraph{予習}
単回帰分析について，数式的表現や最小二乗法の意味について概念的な理解で十分なので再確認しておくこと。重回帰分析では説明変数が複数になるため，解析的算出には行列計算の知識が必要になる。行列計算については，二年次の統計に関する講義で詳しく扱う予定であるが，興味があるものは「線形代数」をキーワードに予習すると良い。初学者には\textcite{Yuki201810}または\textcite{Okada200804}が適している。
\paragraph{復習}
\textcite{Toyoda201706}のP.89--98は本時の内容がわかりやすくまとめられているので，通読しておくと良い復習になる。また，授業で学んだ内容を応用して以下の課題に取り組むこと：

\begin{enumerate}
\item 提供されたデータセットを用いて，Rで重回帰分析を実施する
\item 各説明変数の偏回帰係数と標準化偏回帰係数を求め，その意味を解釈する
\item モデルの決定係数と調整済み決定係数を算出し，モデルの適合度を評価する
\item 多重共線性の診断を行い，問題がある場合は対処法を考察する
\item 交絡要因の影響を統制した結果，単回帰分析と比較してどのような変化があったかを考察する
\end{enumerate}

重回帰分析は心理学研究で広く用いられる手法なので，\textcite{kosugi2016a}の第二章などを読んで，どのような使い方をされているのか，また前提となっている条件や利用上の問題点などについても理解を深めておくこと。
\section{Rをつかった回帰分析}
\label{b09_Regression_onR}

\subsection{授業内容}
\subsubsection{概要}
回帰分析や重回帰分析について，基本的な概念や用語・推定値の意味は前回までに習得済みであるが，それが統計環境では実際どのようにして算出されるか，算出された数字はどのような意味があるかについて具体的に理解することを目的とする。座学の中では記号としてしか意味を持たなかったものが，実際の数字として得られるとより理解が深まることにもなる。統計環境はRを用いるが，ほかの統計パッケージを用いても同様の出力が得られる。統計環境が変わる可能性もあるので，出力画面の位置で内容を覚えるのではなく，意味的な理解をしておく必要がある。
\subsubsection{コマ主題細目}
\begin{description}

	\item[散布図の描画] 線形回帰はその名の通り線形関係をモデルで表現するものであり，非線形な関係や外れ値の存在などがあれば適切な推定値が得られたことにはならない。分析をする前にまずデータの様子を確認するために，散布図の描画は必ず行わなければならない。Rの基本的なplot関数を用いた散布図の描画方法を学び，データの視覚的な確認方法を理解する。また，複数の変数がある場合の散布図行列(pairs関数)の活用方法についても理解する。
	      \begin{flushright}
		      $\to$ \textcite{kosugi2019}のP.92--98.
	      \end{flushright}

	\item[線形回帰の実行] データが読み込まれれば，$lm$関数によって最小二乗法による回帰分析を行うことができる。Rによる関数関係の記述(formula)の方法を理解し，出力結果から適切な読み取りができるよう，どのような用語でどのような出力がなされているかをしっかり確認する。特に，切片(Intercept)と傾き(係数)の推定値，標準誤差，t値，p値の意味と解釈について理解する。
	      \begin{flushright}
		      $\to$ \textcite{kosugi2018}のP.54--60.
	      \end{flushright}

	\item[回帰診断とQ-Qプロット] lm関数によって得られた回帰モデルの適切性を診断することが重要である。回帰分析の結果オブジェクトをplot関数に渡すことで，残差プロット，Q-Qプロット，スケールロケーションプロットなど，モデル診断に有用なグラフが得られる。特にQ-Qプロットは残差の正規性を確認するための重要なツールであり，理論上の正規分布と実際の残差の分布を比較することで，回帰モデルの前提条件が満たされているかを確認できる。残差プロットでは残差のパターン(漏斗状や系統的な曲線など)を確認し，モデルの妥当性や外れ値の影響を検討する方法を学ぶ。
	      \begin{flushright}
		      $\to$ \textcite{kosugi2018}のP.68--72
	      \end{flushright}

	\item[残差と予測値の確認] $lm$関数から返される結果オブジェクトの中には，予測値$\hat{Y}$がfitted.values，残差がresidualsという変数名で保存されている。これらを取り出したり図示したりすることで，回帰分析の諸特徴を理解する一助となる。数理的な導出は\textcite{kosugi2018}のP.127--135にあるが，具体的な数字としてこれを確認し，理解する。また，residuals.lm()，fitted.lm()などの関数を用いた残差や予測値の取り出し方も学ぶ。
	      \begin{flushright}
		      $\to$ \textcite{kosugi2018}のP.54--60.
	      \end{flushright}

	\item[重回帰分析の実行] 統計環境において重回帰分析を行うのは，説明項を1つ追加するだけで良いので作業としては容易い。ただし解釈に注意が必要であることを再確認し，また標準化偏回帰係数の算出方法についても理解を進める。また変数を増やす前と増やした後で，重相関係数が増加することを確認する。さらに，summary関数から得られる調整済み決定係数(Adjusted R-squared)や，AIC，BICなどのモデル選択指標についても理解する。
	      \begin{flushright}
		      $\to$ \textcite{kosugi2018}のP.63--67.
	      \end{flushright}

\end{description}
\subsubsection{キーワード}
\begin{itemize}
	\item Rの基本plot関数
	\item lm関数とformula表記
	\item Q-Qプロットと残差診断
	\item fitted.valueとresiduals
	\item 標準化偏回帰係数
\end{itemize}
\subsection{授業情報}
\paragraph{コマの展開方法}
演習
\paragraph{公認心理師標準シラバスにおける位置づけ}
科目番号5 心理学統計法；②統計に関する基礎的な知識；C エクセル，R，SPSS入門(1)，D エクセル，R，SPSS入門(2) に相当
\paragraph{心理統計法標準カリキュラム(日本心理学会)における位置づけ}
第4回 相関と回帰，第13回 重回帰分析 細目3 に相当
\subsubsection{予習・復習課題}
\paragraph{予習}
RやRStudio，プロジェクトによる管理など基本的なことを改めて教示しないので，前回の復習をかねて空いている時間にPCルームで基本的なR/RStudioの挙動について理解しておくと良い。またRの内部でどのように数字やデータが扱われるかについて，\textcite{kosugi2019}のP.24--36などを参考にみておくと良い。さらに，単回帰分析と重回帰分析の概念的な理解を確認しておくこと。
\paragraph{復習}
授業で扱ったデータを用いて，以下の課題に取り組むこと：

\begin{enumerate}
\item 単回帰分析を実行し，結果を解釈する
\item 回帰診断プロット(特にQ-Qプロット)を描画し，モデルの妥当性を検討する
\item 重回帰分析を実行し，単回帰分析との結果の違いを比較する
\item 残差と予測値を取り出し，散布図上にプロットする
\item 標準化偏回帰係数を算出し，変数間の相対的な影響力を評価する
\end{enumerate}


\section{確率と確率分布}
\label{b10_Probability}

\subsection{授業内容}
\subsubsection{概要}
データ分析においては，不確実性を扱うための枠組みが必要である。本コマでは確率の基本的な考え方と確率分布について学ぶ。確率分布は，起こりうる結果とその生起確率を対応させたものであり，統計学の基盤をなす重要な概念である。最も単純なベルヌーイ分布から始め，二項分布へと拡張し，さらに多数の試行を考えると正規分布に近似することを理解する。これらの確率分布の性質と期待値や分散といった特性値，そして標準化の概念を学ぶことで，後の推測統計学の基礎を形成する。

\subsubsection{コマ主題細目}
\begin{description}
	\item[確率と確率分布の基礎] 確率は不確実な事象を数量的に表現するための道具である。コルモゴロフの公理に基づく確率の基本的性質(非負性，加法性，全確率)について理解し，確率変数と確率分布の概念を学ぶ。離散確率変数の場合は確率質量関数，連続確率変数の場合は確率密度関数によって確率分布が表現される。連続変数の場合，一点の確率は定義できず，区間の面積が確率となることに注意する。また，確率分布の特性値として期待値(平均)と分散の概念を理解する。
	      \begin{flushright}
		      $\to$ \textcite{Maeda2017}のP.78--88, \textcite{Kawabata201408}のP.31--35.
	      \end{flushright}

	\item[ベルヌーイ分布と二項分布] 確率分布の中でも最も基本的なものがベルヌーイ分布である。成功/失敗の二値のみをとる試行(ベルヌーイ試行)の結果を表す確率分布で，成功確率$\theta$をパラメータとする。このベルヌーイ試行をn回繰り返したときの成功回数の分布が二項分布であり，$B(n,p)$と表記される。二項分布の期待値は$np$，分散は$np(1-p)$となることを理解する。また，コイン投げやサイコロ振りなどの具体例を通じて，ベルヌーイ試行の独立性や試行回数の増加に伴う分布の変化について考察する。
	      \begin{flushright}
		      $\to$ \textcite{Nojima201903}のP.66-68, \textcite{Yamada2004}のP.74-80.
	      \end{flushright}

	\item[正規分布] 多くの自然現象や社会現象において観測される値は，しばしば正規分布に従うことが知られている。二項分布の試行回数$n$が大きくなるにつれて，その形状は正規分布に近づく(中心極限定理)。正規分布$N(\mu,\sigma)$は，位置パラメータ$\mu$(平均)と尺度パラメータ$\sigma^2$(分散)によって特徴づけられる連続確率分布である。その確率密度関数は左右対称の釣鐘型であり，平均値を中心に標準偏差の距離によって分布の広がりが決まる。正規分布では平均値から$\pm 1$標準偏差の範囲に約68\%，$\pm 2$標準偏差の範囲に約95\%，$\pm 3$標準偏差の範囲に約99.7\%のデータが含まれるという特徴を持つことを理解する。
	      \begin{flushright}
		      $\to$ \textcite{Naganuma2201611}のP.16--40, \textcite{Kawabata201408}のP.36--40.
	      \end{flushright}

	\item[標準正規分布と確率計算] 様々な平均と標準偏差を持つ正規分布を比較するために，標準化という操作が用いられる。任意の正規分布$N(\mu,\sigma)$から得られたデータ$x$を，$(x-\mu)/\sigma$という変換で標準化すると標準正規分布$N(0,1)$に従うようになる。標準正規分布では，確率の計算や比較が容易になる。標準正規分布表やRなどの統計ソフトウェアを用いた確率密度(dnorm)，累積確率(pnorm)，確率点(qnorm)の計算方法を学び，実際のデータ分析における活用法を理解する。また，z値(標準化得点)の意味と解釈，偏差値への変換方法についても理解を深める。
	      \begin{flushright}
		      $\to$ \textcite{Kawabata201408}のP.36--44, \textcite{Yamada2004}のP.86--89.
	      \end{flushright}
\end{description}

\subsubsection{キーワード}
\begin{itemize}
	\item ベルヌーイ分布と二項分布
	\item 確率変数と確率分布
	\item 期待値と分散
	\item 正規分布と標準正規分布
	\item z値と標準化
\end{itemize}

\subsection{授業情報}
\paragraph{コマの展開方法}
講義
\paragraph{公認心理師標準シラバスにおける位置づけ}
科目番号5 心理学統計法；①心理学で用いられる統計手法；D 離散分布：二項分布を中心として，E 連続分布：正規分布を中心として／科目番号5 心理学統計法；②統計に関する基礎的な知識；B 確率と確率分布 に相当
\paragraph{心理統計法標準カリキュラム(日本心理学会)における位置づけ}
第7回 母集団と標本 細目3 に相当
\subsubsection{予習・復習課題}
\paragraph{予習}
高校数学で学んだ確率・統計の内容を復習しておくこと。特に，確率の基本法則(加法定理，乗法定理)，順列・組合せの考え方について確認しておくとよい。また，確率分布について入門的な解説をしている\textcite{Hamada201812}や\textcite{Maeda2017}の最初の数章に目を通しておくと理解の助けになるだろう。
\paragraph{復習}
コイン投げや賽子を振るといった具体的な例を通じて，ベルヌーイ分布と二項分布の性質を確認してみよう。また，正規分布の特徴である「左右対称」「単峰性」について\textcite{Naganuma2201611}を参照しながら理解を深めよう。Rを用いて，様々なパラメータを持つ二項分布と正規分布のグラフを描画し，パラメータの変化に伴う分布の形状の変化を視覚的に確認することも有効である。以下のコマンドを試してみるとよい：

\begin{enumerate}
\item 二項分布のグラフ: `plot(0:20, dbinom(0:20, size=20, prob=0.5), type="h")`
\item 正規分布の確率密度関数: `curve(dnorm(x, mean=0, sd=1), from=-4, to=4)`
\item 標準正規分布の累積分布関数: `curve(pnorm(x), from=-4, to=4)`
\item データの標準化と確率の計算: `pnorm((x-mean(x))/sd(x))`
\end{enumerate}

また，日常生活や心理学研究で見られる様々な現象の中から正規分布に近似できるものを考え，その理由について考察してみよう。
\section{確率と統計的推論}
\label{b11_estimation}

\subsection{授業内容}
\subsubsection{概要}
母集団に確率分布を仮定し，そこから得られる標本分布の特徴を利用して母数を推定する方法について学ぶ。標本の平均値が母平均に一致することを確認し，また標本の分散は母分散に一致しないことを確認する。これまでは手元にあるデータを記述する方法を学んできたが，本コマからは未知の母集団の特性を推測するための統計的推論の基礎に入る。現実の研究では全数調査は困難であるため，標本から母集団へと推論することが必要となる。ここでの推定方法はモーメントを用いたものであるが，確率分布を用いる方法もあり，その場合の``確率モデル''の考え方が必要なことを学ぶ。

\subsubsection{コマ主題細目}
\begin{description}
	\item[母集団と標本] 推測統計学の基本的な概念である母集団と標本の関係を理解し，統計学的な推定(点推定，区間推定)の位置付けについて学ぶ。\footnote{本書では検定をモデル比較の一種として扱うので，ここでは言及しない。}用語としての母数(母平均，母分散，母比率など)と標本統計量(標本平均，標本分散，標本比率など)，および標本統計量の実現値といった用語の意味するところも慎重に理解しなければならない。特に母集団分布，標本分布，サンプリング分布の違いについて理解することが重要である。
	      \begin{flushright}
		      $\to$ \textcite{Yamada2004}のP.68--73
	      \end{flushright}

	\item[母集団分布とモデル] ここでは母集団に正規分布を仮定する``モデル''について考える。統計的な推測は基本的に不良設定問題であり，何らかの前提をおかないと解くことができない問題がほとんどである。心理学の場合は母集団分布を正規分布とする仮定が選ばれることが多く，この場合は標本統計量の分布も計算しやすくなる。現実の分布が必ずしも正規分布に従わない場合でも，中心極限定理により標本平均は正規分布に近似できることが，この仮定の有用性を支えている。
	      \begin{flushright}
		      $\to$ \textcite{Yamada2004}のP.74--79.
	      \end{flushright}

	\item[正規分布の別の意味] 正規分布はここまで，誤差の分布として解説してきたが，複数の要因が積み重なった時の全体的傾向として理解することもできる。平均値が集団の代表値としての意味もあること，また中心極限定理によって分布によらずその標本平均値が正規分布に従うことが示される。中心極限定理は，母集団分布がどのような形であっても，標本サイズが十分大きければ標本平均の分布は正規分布に近づくという重要な性質であり，推測統計学の基礎となる定理である。
	      \begin{flushright}
		      中心極限定理については$\to$ \textcite{Minamoto2015}のP.180--181.
	      \end{flushright}

	\item[標本分布と標準誤差] 母集団が正規分布に従うとすると，標本統計量も正規分布に従うことが導出できる。ここでは標本統計量としての平均値を例にあげ，「標本統計量の分布」のことを標本分布と呼ぶこと，また標本分布の標準偏差をとくに標準誤差と呼ぶことを理解する。標準誤差は標本平均のばらつきを表す指標であり，サンプルサイズの平方根に反比例することを理解する。たとえばサンプルサイズが大きくなったら，小さくなったらどのような値が得られるのかについて，具体的な例をみながら考えると理解が進む。
	      \begin{flushright}
		      $\to$ 標本分布については\textcite{Yamada2004}のP.90--97.
	      \end{flushright}

	\item[標準正規分布と統計的推論] 標準正規分布は統計的推論において中心的な役割を果たす。標本統計量を標準化することで標準正規分布に対応させ，その確率的な性質を利用して推定や後の検定を行うことができる。標準正規分布表の使い方や，標準正規分布に基づく信頼区間の構成方法について学ぶ。また，標準正規分布を用いた統計的推論の基本的な考え方と，これが後の検定につながることを理解する。
	      \begin{flushright}
		      $\to$ \textcite{Yamada2004}のP.86--89，\textcite{Kawabata201408}のP.36--44.
	      \end{flushright}
\end{description}
\subsubsection{キーワード}
\begin{itemize}
	\item 母集団と標本
	\item 標本分布と標本統計量
	\item 中心極限定理
	\item 標準誤差
\end{itemize}
\subsection{授業情報}
\paragraph{コマの展開方法}
講義
\paragraph{公認心理師標準シラバスにおける位置づけ}
科目番号5 心理学統計法；②統計に関する基礎的な知識；A 統計分析の基礎：母集団と標本 に相当
\paragraph{心理統計法標準カリキュラム(日本心理学会)における位置づけ}
第7回 母集団と標本 細目1・2，第8回 統計的推定 細目2 に相当
\subsubsection{予習・復習課題}
\paragraph{予習}
確率と確率分布の基本概念を復習しておくこと。特に標準正規分布の特徴とその利用方法について確認しておくと理解が深まる。また，ここでは\textcite{Yamada2004}のP.68-97が丁寧に解説されているので，事前に一読しておくことを強く勧める。
\paragraph{復習}
いよいよ記述統計をこえ，直接知り得ないものについて推論する段階に進む。基本的に解けない問題を解こうとするようなものであり，仮定やモデル，みなしている箇所について自覚的でなければならない。今回の講義の中で，何が仮定され，何が理論的に導出されたものなのかを区別できるかどうか，ノートを整理しておくと良いだろう。また推測統計学の基本的概念である，母集団と標本の関係と推定の手続きについては，\textcite{Yamada2004}のテキストが最も丁寧で適切な記述をしているので，是非該当箇所を通読し，概念や用語を間違えないように注意されたし。

以下の点について理解を深めておくことも有効である：
\begin{enumerate}
\item 標本サイズと標準誤差の関係
\item 中心極限定理の意味とその応用
\item 母集団分布を正規分布と仮定することの意義と限界
\item 点推定と区間推定の違いと用途
\end{enumerate}
\section{推定}
\label{b12_test}

\subsection{授業内容}
\subsubsection{概要}
これまでの授業では母集団分布と標本分布について学んできた。本コマでは標本から母集団の特性値(母数)を推定する方法について学ぶ。まず不偏推定量の概念を理解し，点推定と区間推定の違いについて学ぶ。標本統計量を用いて母数を推定する際の基本的な考え方と，推定値の精度や信頼性を評価する方法を理解する。さらに，推定と検定の関係についても触れ，統計的推論の基本的な枠組みを身につける。また，標準正規分布が統計的推論において果たす重要な役割についても理解を深める。

\subsubsection{コマ主題細目}
\begin{description}
	\item[不偏推定量] 母集団に正規分布を仮定し，標本平均の平均(期待値)が母平均に一致することを確認する。このことから，標本統計量が推定値として利用できることを確認する。ただし，分散の場合は一致しない(不偏性がない)ことをみる。分散については不偏分散を考えることでこの不一致を解決することができる。標本統計量が母数の良い推定量であるための条件(不偏性，一致性，効率性)についても理解する。また，推定量と推定値の違いについても理解する。
	      \begin{flushright}
		      $\to$ 不偏推定量については\textcite{Yamada2004}のP.98--103に詳しい。標本統計量のこれらの特徴については，\textcite{kosugi2023}のP.119--141も参照すると良い。
	      \end{flushright}

	\item[点推定と標準誤差] 点推定とは，標本統計量を用いて母数の値を単一の値として推定する方法である。例えば，標本平均は母平均の点推定値となる。しかし，この推定値にはサンプリング誤差が含まれるため，その精度を評価する必要がある。標準誤差は推定値の精度を表す重要な指標であり，標本平均の標準誤差は母標準偏差をサンプルサイズの平方根で割ったものとして定義される。標準誤差の意味と解釈，およびサンプルサイズとの関係について理解する。また，中心極限定理により，サンプルサイズが大きくなると標本平均の分布が正規分布に近づくことも確認する。
	      \begin{flushright}
		      $\to$ \textcite{Kawabata201408}のP.107--110，\textcite{Yamada2004}のP.116--119.
	      \end{flushright}

	\item[区間推定と信頼区間] 標本平均は母平均の不偏推定量ではあるが，ある値をそのまま母平均であると考えるのは(標準誤差の考え方があるとは言え)やや断定的である。これに対して一定の幅を使って予測する区間推定という方法がある。推定に確率の言葉が入ってくるので不慣れな表現が出てくるが，注意深く解釈する必要がある。95\%信頼区間の意味と解釈，標準誤差を用いた信頼区間の計算方法について理解する。また，信頼区間の幅とサンプルサイズの関係についても学ぶ。信頼区間の正しい解釈は，「この方法で区間を構成すると，長期的には95\%の確率で真の母数を含む」というものであり，特定の区間が母数を含む確率が95\%という意味ではないことに注意する。
	      \begin{flushright}
		      $\to$ \textcite{Kawabata201408}のP.111--116，\textcite{Yamada2004}のP.120--127.
	      \end{flushright}


	\item[推定と検定の関係] 推定は母数の値を推し量る作業であるのに対し，検定はある仮説の下での統計量の期待値と実際の統計量を比較することで，仮説の妥当性を判断する作業である。両者は密接に関連しており，推定は「どれくらいか」という問いに，検定は「差があるか否か」という問いに答えるものである。推定と検定の関係を理解することで，次回学ぶ帰無仮説検定の基礎を形成する。特に，信頼区間と検定の関係について理解を深める。
	      \begin{flushright}
		      $\to$ \textcite{Kawabata201408}のP.81--85，\textcite{Yamada2004}のP.104--110.
	      \end{flushright}

\end{description}
\subsubsection{キーワード}
\begin{itemize}
	\item 不偏推定量
	\item 点推定と標準誤差
	\item 区間推定と信頼区間
	\item 標準正規分布
	\item 推定と検定の関係
\end{itemize}
\subsection{授業情報}
\paragraph{コマの展開方法}
講義
\paragraph{公認心理師標準シラバスにおける位置づけ}
科目番号5 心理学統計法；①心理学で用いられる統計手法；F 推測統計：その考え方 に相当
\paragraph{心理統計法標準カリキュラム(日本心理学会)における位置づけ}
第8回 統計的推定 細目1・2・3・4 に相当
\subsubsection{予習・復習課題}
\paragraph{予習}
前回までに学んだ標本分布，標準誤差，標準正規分布についての理解を確認しておくこと。特に，標本統計量の標準化の方法と標準正規分布表の読み方について復習しておくと理解が深まる。
\paragraph{復習}
授業で学んだ内容を定着させるために，以下の点について理解を深めておくこと：

\begin{enumerate}
\item 不偏推定量の概念と，標本分散と不偏分散の違い
\item 点推定と標準誤差の関係，およびサンプルサイズが推定精度に与える影響
\item 信頼区間の正しい解釈と，その構成方法
\item 標準正規分布を用いた統計的推論の基本的な考え方
\item 推定と検定の関係，および信頼区間と検定の関連性
\end{enumerate}

統計的推定は，実証的な心理学研究の基盤をなす重要な概念である。特に，点推定だけでなく区間推定の考え方を理解することで，データの持つ不確実性を適切に評価できるようになる。統計的推定についての理解を深めるには，\textcite{Kawabata201408}のP.104--116や\textcite{Yamada2004}のP.98--127などを参考にすると良い。
\section{検定}
\label{b13_NHST}

\subsection{授業内容}
\subsubsection{概要}
前回は統計的推定について学んだ。本コマでは，推定を発展させた統計的意思決定法である帰無仮説検定について学ぶ。帰無仮説検定は心理学研究において広く用いられている統計的推論の方法であり，データに基づいて仮説の妥当性を評価するための枠組みを提供する。本コマでは，帰無仮説検定の基本的な考え方と手順，さまざまな検定法(相関係数の検定，カイ二乗検定など)，そして有意水準や片側・両側検定の概念について理解を深める。帰無仮説検定は背理法を用いて効果の有無を判断する方法であり，結果の報告に関する表記法も含めてその作法を学ぶ必要がある。

\subsubsection{コマ主題細目}
\begin{description}
	\item[統計的帰無仮説検定とその手順] 統計的帰無仮説検定(Null Hypothesis Significance Test; NHST)は帰無仮説に基づいてモデル比較をする一連の手続きである。前回学んだ推定と異なり，検定では特定の仮説(帰無仮説)が正しいと仮定した場合に観測されたデータがどれだけ起こりにくいかを評価する。帰無仮説検定の基本的な考え方，帰無仮説と対立仮説の設定，検定統計量の算出，p値の計算，判断基準の設定という一連の流れを学ぶ。また，検定統計量は標本統計量を標準化したものであり，統計量によって適切な分布(標準正規分布，t分布，カイ二乗分布，F分布など)が対応することを理解する。
	      \begin{flushright}
		      $\to$ 手続きの一般化としては\textcite{Yamada2004}のP.108--109，\textcite{Kawabata201408}のP.81--85.
	      \end{flushright}

	\item[相関係数の検定] 相関係数の統計的検定は，母相関係数がゼロであるという帰無仮説を検証するものである。2変数が無相関(母相関係数$\rho=0$)という帰無仮説の下では，サンプルサイズ$n$と標本相関係数$r$から計算される検定統計量が自由度$n-2$の$t$分布に従うことを理解する。この検定統計量の計算方法と，検定結果の解釈について学ぶ。また，相関係数の検定と同時に，相関係数の信頼区間を求める方法についても触れる。さらに，相関係数の大きさと統計的有意性の違いについても理解を深める。無相関検定はRでは`cor.test()`関数を用いて実行できる。
	      \begin{flushright}
		      $\to$ \textcite{Kawabata201408}のP.98--103，\textcite{kosugi2018}のP.58--60，\textcite{Shimizu2021}のP.150--162.
	      \end{flushright}

	\item[カイ二乗検定] カイ二乗検定は，カテゴリカルデータの分析に用いられる検定法である。観測度数と期待度数の差を評価し，分布の適合度や変数間の独立性を検証する。カイ二乗検定には適合度検定と独立性の検定があり，それぞれの適用場面と計算方法について学ぶ。カイ二乗統計量の算出方法と，カイ二乗分布を用いた確率評価の方法について理解する。また，カイ二乗検定の適用条件(期待度数が5以上など)や結果の解釈についても学ぶ。
	      \begin{flushright}
		      $\to$ カイ二乗検定については\textcite{Kawabata201408}のP.145--153，\textcite{Yamada2004}のP.177--183を参照。
	      \end{flushright}

	\item[1つの平均値の検定(t検定)] 前回は相関係数の検定とカイ二乗検定について学んだが，本コマでは平均値の検定について学ぶ。実際の研究では母分散がわかっていることはほとんどなく，標本から不偏分散を用いて推定する必要がある。このような場合は標準正規分布ではなく，$t$分布と呼ばれる統計量が用いられる。$t$分布の特徴(自由度によって形状が変わる，サンプルサイズが大きくなると標準正規分布に近づくなど)，$t$検定の手順，および$t$分布表の使い方について学ぶ。また，1サンプルの$t$検定と，次回学ぶ2サンプルの$t$検定の違いについても理解する。
	      \begin{flushright}
		      $\to$ \textcite{Kawabata201408}のP.93--96，\textcite{Yamada2004}のP.128--131.
	      \end{flushright}
\end{description}
\subsubsection{キーワード}
\begin{itemize}
	\item 帰無仮説検定と検定統計量
	\item 相関係数の検定(無相関検定)
	\item カイ二乗検定
	\item 1つの平均値の検定(t検定)
\end{itemize}
\subsection{授業情報}
\paragraph{コマの展開方法}
適宜投影資料を用いる
\paragraph{公認心理師標準シラバスにおける位置づけ}
科目番号5 心理学統計法；①心理学で用いられる統計手法；F 推測統計：その考え方，I 推測統計：頻度や比率をめぐって に相当
\paragraph{心理統計法標準カリキュラム(日本心理学会)における位置づけ}
第9回 統計的仮説検定 細目1・3・4 に相当
\subsubsection{予習・復習課題}
\paragraph{予習}
前回学んだ不偏推定量，点推定と区間推定，標準正規分布の知識を復習しておくこと。特に，標準誤差と信頼区間の概念，および推定と検定の関係について確認しておくと理解が深まる。また，相関係数の概念と計算方法についても再確認しておくこと。
\paragraph{復習}
授業で学んだ内容を定着させるために，以下の点について理解を深めておくこと：

\begin{enumerate}
\item 帰無仮説検定の論理構造と基本的な手順
\item 相関係数の検定の手順と解釈(無相関の検定，検定統計量の計算，結果の解釈)
\item カイ二乗検定の適用場面と使い方(適合度検定と独立性の検定)
\item 有意水準とp値の関係，およびp値の正しい解釈
\item 片側検定と両側検定の違いと，それぞれの適用場面
\end{enumerate}

また，実際に検定を行ってみることで理解が深まる。例えば，相関係数の大きさ，標本の大きさを任意の値にし，無相関検定を実際にやってみると良い。さまざまな数値を入れることで，検定結果がどのように変動するかを理解することは帰無仮説検定の本質を理解するのに役立つ。仮想的な数字では実感が得られにくい場合，身の回りの相関がありそうなデータを具体例として検定し，何が言えるのかを考えてみると良い。
\section{帰無仮説検定の応用と注意点}
\label{b14_NHST2}

\subsection{授業内容}
\subsubsection{概要}
前回までに帰無仮説検定の基本的な考え方と手順について学んだ。本コマでは，より実践的な検定法として母分散が未知の場合の平均値の検定(t検定)を学び，検定結果の報告方法や解釈について理解を深める。また，帰無仮説検定における2種類の誤り(第1種の誤りと第2種の誤り)について学び，検定力と効果量の概念を理解する。さらに，帰無仮説検定の限界や注意点についても考察する。特に，統計的有意性と実質的重要性の違い，検定の多重性の問題などについて理解し，帰無仮説検定を適切に活用するための視点を身につける。

\subsubsection{コマ主題細目}
\begin{description}

	\item[有意水準と$p$値の解釈]
	モデル比較という観点からは，どうしても「判定基準」を考える必要がある。NHSTの文脈では，それは希少性の指標である有意水準になる。有意水準の定め方は領域ごとに定める任意であり，心理学では一般に5\%が用いられている。このことは，複数のモデルの優劣を判定する数字であって，モデルの強さ，正しさを表現する数字ではないことに注意する。有意水準とp値の関係，および有意水準の設定の仕方について理解する。また，p値の正しい解釈についても学び，p値の誤解や誤用についても理解を深める。
	      \begin{flushright}
		      $\to$ \textcite{Yamada2004}のP.112--115，\textcite{Kawabata201408}のP.86--89.
	      \end{flushright}
	      \begin{flushright}
		      $\to$ 検定の多重性については\textcite{Yamada2004}のP.158--161，帰無仮説検定の限界については\textcite{Kawabata201408}のP.99--103.
	      \end{flushright}
	
		  
	\item[結果の報告と解釈] 帰無仮説検定の結果を論文やレポートで報告する際の適切な方法について学ぶ。統計量の値，自由度，p値，効果量などを正確に報告することの重要性を理解する。また，p値の捉え方として「帰無仮説が正しいという仮定のもとで，観測されたデータやそれよりも極端なデータが得られる確率」であることを再確認し，p値の誤った解釈(帰無仮説が正しい確率，効果の大きさなど)を避けるようにする。さらに，統計的に有意であることの意味と限界について理解を深める。
	      \begin{flushright}
		      $\to$ \textcite{Yamada2004}のP.132--133，\textcite{Yamauchi201001}のP.216--220.
	      \end{flushright}

	\item[検定における2種類の誤り] 帰無仮説検定は確率的な判断を行う手法であるため，誤った判断を下す可能性が常に存在する。検定における誤りには2種類あり，第1種の誤り(Type I Error)は帰無仮説が正しいのに棄却してしまう誤り，第2種の誤り(Type II Error)は帰無仮説が誤りなのに棄却できない誤りである。有意水準αは第1種の誤りを犯す確率の上限を表し，検定力(1-β)は第2種の誤りを犯さない確率を表す。検定力はサンプルサイズ，効果量，有意水準によって決まることを理解し，適切なサンプルサイズの決定方法(検出力分析)についても学ぶ。
	      \begin{flushright}
		      $\to$ \textcite{Yamada2004}のP.120--121，検定の多重性については同書P.158--161，\textcite{Kawabata201408}の P.90--93も参照。
	      \end{flushright}

	\item[効果量と例数設計] 統計的有意性(p値が有意水準より小さいこと)は必ずしも実質的な重要性を意味するわけではない。特に大きなサンプルサイズでは，わずかな効果でも統計的に有意になりやすい。そこで，効果の大きさを標準化して表す効果量の概念が重要になる。代表的な効果量として，相関係数(r)，Cohen's d，η²(イータ二乗)などがあり，それぞれの意味と解釈について学ぶ。また，効果量の大きさを評価する目安(小・中・大)についても理解する。効果量は検定結果とともに報告することで，結果の実質的な意味を評価する助けとなる。
	      \begin{flushright}
		      $\to$ 効果量については\textcite{Matia201201}のP.43--46，\textcite{Kawabata201408}のP.97--98.
	      \end{flushright}




\end{description}
\subsubsection{キーワード}
\begin{itemize}
	\item 有意水準とp値
	\item 片側検定と両側検定
	\item 第1種の誤りと第2種の誤り
	\item 検定力と効果量
	\item 統計的有意性と実質的重要性
	\item 検定の多重性と帰無仮説検定の限界
\end{itemize}
\subsection{授業情報}
\paragraph{コマの展開方法}
適宜投影資料を用いる
\paragraph{公認心理師標準シラバスにおける位置づけ}
科目番号5 心理学統計法；①心理学で用いられる統計手法；F 推測統計：その考え方／科目番号4 心理学研究法；②データを用いた実証的な思考方法；C データの統計的記述 に相当
\paragraph{心理統計法標準カリキュラム(日本心理学会)における位置づけ}
第9回 統計的仮説検定 細目2，第12回 統計的仮説検定の注意点 細目1 に相当
\subsubsection{予習・復習課題}
\paragraph{予習}
前回学んだ帰無仮説検定の基本的な考え方や手順，有意水準とp値の関係について復習しておくこと。特に，相関係数の検定やカイ二乗検定の手順と解釈について確認しておくと理解が深まる。
\paragraph{復習}
授業で学んだ内容を定着させるために，以下の点について理解を深めておくこと：

\begin{enumerate}
\item t検定の手順と適用条件，およびt分布の特徴
\item 第1種の誤りと第2種の誤りの違い，および検定力との関係
\item 効果量の種類と解釈，および統計的有意性との違い
\item 帰無仮説検定の結果を適切に報告する方法
\item 帰無仮説検定の限界と注意点，および対処法
\end{enumerate}

検定の話に入ると，要因計画や線形モデルとどのようにこれらが接合するのかが分かりにくい。回帰モデルの説明変数を離散化したのが要因計画→要因計画の効果の大きさ(傾きの大きさ)を質的に判断する手続きが，平均値差の検定である，という関係であることを整理しておこう。推測統計学は，基本的に母数の推定(回帰係数のパラメータも母数である)をしているのであり，推定量を何らかの基準で判断するのが検定，という繋がりである。もちろん検定を行わずに，大きさをそのまま報告しても良い。むしろ昨今ではそちらの方が推奨されており，些細な違いや「差がないこと」を「差があること」の根拠とすることの問題点が数多く指摘されている。帰無仮説検定の手続きは，過去の論文に示された結果を読むための知識として重要であり，本質である効果の大きさを推定していることを忘れないように。
\section{前期のまとめと試験}
\label{b15_summary_exam}

\subsection{授業内容}
\subsubsection{概要}
これまでの14回の授業を通じて，心理統計の基礎となる概念と手法について学んできた。本コマでは前期の内容を総括するとともに，学習内容の理解度を確認するための試験を実施する。試験は授業時間内に行われ，これまでに学んだ記述統計，相関と回帰，確率分布，推測統計と検定などの内容から出題される。試験を通じて，心理統計の基本概念の理解と実践的な応用力が身についているかを確認する。

\subsubsection{コマ主題細目}
\begin{description}
	\item[前期の総括] 前期では「記述」と「因果と相関」を中心に学んできた。心理学における数量的アプローチの意義から始まり，データの可視化，記述統計量，相関と回帰分析，確率分布と統計的推論，そして帰無仮説検定まで，心理統計の基礎となる概念と手法を体系的に学習してきた。これらの知識は，後期の応用的な内容(分散分析，因子分析など)を学ぶための基盤となる。また，Rを用いた実習を通じて，統計解析の実践的スキルも身につけてきた。前期で学んだ内容を振り返り，後期の学習へとつなげる。
	      \begin{flushright}
		      $\to$ これまでの各回の内容を復習しておくこと。
	      \end{flushright}

	\item[学習到達度の確認試験] 前期で学んだ内容の理解度を確認するため，授業時間内に試験を実施する。試験はマークシート形式で行われ，A4用紙2枚までの自筆メモの持ち込みが許可される。テスト理論に基づいて等価性が担保された問題が出題され，前期の学習内容全般から幅広く出題される。特に，記述統計の基本概念，相関と回帰分析の解釈，確率分布の特性，推測統計と帰無仮説検定の考え方など，心理統計の基本的な理解が問われる。
	      \begin{flushright}
		      $\to$ 試験の詳細については，別途配布される「試験に関するお知らせ」を参照のこと。
	      \end{flushright}

	\item[後期の学習に向けて] 前期では主にデータの記述と基本的な統計的推論について学んだ。後期ではより複雑な統計モデルや分析手法(分散分析，多変量解析など)について学び，実際の心理学研究に応用する方法を学ぶ。前期で習得した基礎知識をもとに，より高度な統計的思考と分析技術を身につけていく。特に，サンプルから母集団への一般化や，複数の要因が絡む実験計画の分析方法など，心理学研究の実践に直結する内容が中心となる。
	      \begin{flushright}
		      $\to$ 夏休み中に前期の内容を復習し，後期の学習に備えること。
	      \end{flushright}

\end{description}
\subsubsection{キーワード}
\begin{itemize}
	\item 記述統計と推測統計
	\item 相関と因果
	\item 確率分布と統計的推論
	\item 帰無仮説検定
	\item 統計的思考の実践
\end{itemize}
\subsection{授業情報}
\paragraph{コマの展開方法}
前期内容の総括(約15分)の後，授業時間内試験(約60分)を実施する
\paragraph{公認心理師標準シラバスにおける位置づけ}
科目番号5 心理学統計法の総括
\paragraph{心理統計法標準カリキュラム(日本心理学会)における位置づけ}
総括
\subsubsection{予習・復習課題}
\paragraph{予習}
前期の全授業内容を復習しておくこと。特に以下の点について理解を深めておくことが重要である：
\begin{enumerate}
\item 記述統計の基本(平均，分散，標準偏差など)とその意味
\item 相関係数と回帰分析の基本概念と解釈
\item 確率分布(二項分布，正規分布など)の特徴と用途
\item 標本分布と中心極限定理の考え方
\item 推測統計の基本概念(点推定，区間推定)
\item 帰無仮説検定の論理と手順
\end{enumerate}
また，授業で扱ったRの基本操作とデータ分析の方法についても振り返っておくこと。試験では自筆のメモ(A4用紙2枚まで，両面可)の持ち込みが許可されるので，重要な概念や公式をまとめておくと良い。

\paragraph{復習}
試験終了後は，出題された問題について再度考え，不明点があれば教科書や参考書，授業ノートで確認しておくこと。また，後期の学習に向けて，以下の点を意識しておくと良い：

\begin{enumerate}
\item 心理学研究における統計的思考の重要性
\item データに基づく客観的な推論の方法
\item 統計ソフトウェア(R)を用いた実践的な分析スキル
\item 心理学の諸問題を統計的に捉える視点
\end{enumerate}

夏休み期間中に前期の内容を総復習し，特に苦手な分野があれば重点的に学習しておくことを推奨する。後期ではより高度な統計手法を学ぶため，前期の基礎的な内容をしっかりと理解しておくことが重要である。
\section{Rを使ったt検定と検定力分析}
\label{b16_ttest_R}

\subsection{授業内容}
\subsubsection{概要}
夏休みを挟んで，前期に学んだ帰無仮説検定の基本的な考え方を復習し，統計ソフトウェアRを用いた実践的な分析方法を学ぶ。特に，心理学研究でよく用いられるt検定(対応のあるt検定と対応のないt検定)に焦点を当て，Rによる実行方法とその結果の解釈について理解を深める。さらに，シミュレーションを通じて効果量とサンプルサイズが検定結果に与える影響を体験的に学び，適切な研究デザインの重要性について理解する。また，検定力分析の基本的な考え方と実行方法についても学ぶ。

\subsubsection{コマ主題細目}
\begin{description}
	\item[帰無仮説検定の復習] 前期に学んだ帰無仮説検定の基本的な考え方と手順について復習する。帰無仮説と対立仮説の設定，検定統計量の算出，p値の解釈，有意水準に基づく判断という一連の流れを確認する。特に，検定結果の解釈と報告方法について理解を深め，統計的有意性と実質的重要性の違いについても再確認する。また，t分布の特徴と，標準正規分布との関係についても理解する。
	      \begin{flushright}
		      $\to$ \textcite{Kawabata201408}のP.81--96，\textcite{Yamada2004}のP.104--119.
	      \end{flushright}

	\item[対応のあるt検定] 対応のあるt検定(一要因被験者内計画の2水準間の比較)の考え方と実行方法について学ぶ。この検定は同一被験者から得られた2つの条件下でのデータ(事前・事後測定など)を比較する場合に用いられる。Rでの実行方法として，`t.test()`関数の使い方，特に`paired=TRUE`オプションの意味と使用方法を理解する。また，結果の読み方と解釈，効果量の算出方法についても学ぶ。実際のデータセットを使用した演習を通じて，対応のあるt検定の実践的な実行方法を身につける。
	      \begin{flushright}
		      $\to$ \textcite{kosugi2019}のP.126--132，\textcite{Shimizu2021}のP.124--130.
	      \end{flushright}

	\item[対応のないt検定] 対応のないt検定(独立した2群の平均値の比較)の考え方と実行方法について学ぶ。この検定は異なる被験者群から得られたデータを比較する場合に用いられる。Rでの実行方法として，`t.test()`関数の基本的な使い方と，分散の等質性を仮定する場合(`var.equal=TRUE`)と仮定しない場合(Welchの補正)の違いについて理解する。また，検定結果の読み方と解釈，効果量の算出方法についても学ぶ。実際のデータセットを使用した演習を通じて，対応のないt検定の実践的な実行方法を身につける。
	      \begin{flushright}
		      $\to$ \textcite{kosugi2019}のP.132--138，\textcite{Shimizu2021}のP.130--136.
	      \end{flushright}

	% \item[効果量とサンプルサイズ] Rを用いたシミュレーションを通じて，効果量とサンプルサイズが検定結果に与える影響について体験的に学ぶ。異なる効果量(小・中・大)とサンプルサイズの組み合わせでデータを生成し，t検定を実行することで，どのような条件で統計的有意性が得られやすいか，または得られにくいかを確認する。これにより，効果量が小さい場合は大きなサンプルサイズが必要であること，サンプルサイズが小さい場合は検出力が低下することなどを理解する。また，シミュレーションの結果から，標本サイズを適切に設定することの重要性についても学ぶ。
	%       \begin{flushright}
	% 	      $\to$ 効果量については\textcite{Matia201201}のP.43--46，シミュレーションの基本的な考え方については\textcite{Kinosady2021}のP.178--186.
	%       \end{flushright}

	\item[検定力分析の基礎] 検定力(第2種の誤りを犯さない確率)の概念と，研究デザインにおける重要性について学ぶ。特に，Rの`power.t.test()`関数を用いた検定力分析の実行方法と解釈について理解する。事前の検定力分析(研究計画段階での必要サンプルサイズの算出)と事後の検定力分析(研究結果の評価)の違いと用途について学び，研究デザインの質を高めるための手法として検定力分析を活用する方法を身につける。また，サンプルサイズ，効果量，有意水準，検定力の関係について理解を深める。
	      \begin{flushright}
		      $\to$ 検定力分析については\textcite{Kawabata201408}のP.90--93，\textcite{kosugi2023}のP.187--224も参照。
	      \end{flushright}

\end{description}
\subsubsection{キーワード}
\begin{itemize}
	\item 対応のあるt検定と対応のないt検定
	\item Rによる統計分析の実行
	\item 効果量とCohen's d
	\item サンプルサイズとシミュレーション
	\item 検定力分析
\end{itemize}
\subsection{授業情報}
\paragraph{コマの展開方法}
講義と演習
\paragraph{公認心理師標準シラバスにおける位置づけ}
科目番号5 心理学統計法；①心理学で用いられる統計手法；G 推測統計：代表値と散布度をめぐって(1) に相当
\paragraph{心理統計法標準カリキュラム(日本心理学会)における位置づけ}
第10回 平均値の差の検定 細目1・2・3・4・5 に相当
\subsubsection{予習・復習課題}
\paragraph{予習}
前期に学んだ帰無仮説検定の基本的な考え方と手順について復習しておくこと。特に，帰無仮説と対立仮説の設定，検定統計量の算出，p値の解釈，有意水準に基づく判断という一連の流れを確認しておくとよい。また，RとRStudioの基本的な操作方法についても再確認しておくこと。前期に作成したスクリプトやRmarkdownファイルを見直し，基本的なコマンドや関数の使い方を思い出しておくことが望ましい。

\paragraph{復習}
授業で学んだ内容を定着させるために，以下の課題に取り組むこと：

\begin{enumerate}
\item 提供されたデータセットを用いて，対応のあるt検定と対応のないt検定をそれぞれ実行し，結果をRmarkdownで適切にレポートにまとめる
\item 異なる効果量とサンプルサイズの組み合わせでシミュレーションデータを生成し，t検定を実行して結果を比較する
\item 与えられた研究シナリオに対して，適切なサンプルサイズを算出するための検定力分析を実行する
\item 検定結果の統計的有意性と効果量の大きさの両方を考慮した解釈を行い，研究における意義について考察する
\end{enumerate}

これらの課題をRmarkdownにまとめ，提出すること。また，検定力分析とサンプルサイズの重要性について理解を深めるために，\textcite{Matia201201}や\textcite{Kawabata201408}，\textcite{kosugi2023}の関連する章を読むことが推奨される。
\section{実験計画法の基礎とRCT}
\label{b17_design_RCT}

\subsection{授業内容}
\subsubsection{概要}
心理学研究では，変数間の因果関係を検証するために実験的アプローチが重要な役割を果たす。本コマでは，実験計画法の基礎的概念と，特にランダム化比較試験(Randomized Controlled Trial; RCT)の意義と方法について学ぶ。RCTは因果関係を検証するための強力な方法であり，その理論的背景と実践的手法を理解することは，心理学研究の質を高めるために不可欠である。また，実験における誤差の概念とその分布特性についても学び，統計的推論との関連を理解する。実験計画と統計的分析の橋渡しとなる基本的な考え方を身につけることを目指す。

\subsubsection{コマ主題細目}
\begin{description}
	% \item[実験計画法の基本的考え方] 心理学における実験的アプローチの意義と限界について理解する。実験とは，操作した独立変数が従属変数に与える効果を測定することで因果関係を検証する方法である。相関研究と実験研究の違い，因果関係を主張するための条件(時間的先行性，共変関係，他の説明可能性の排除)について学ぶ。また，実験の内的妥当性と外的妥当性のバランス，実験における倫理的配慮についても理解する。
	%       \begin{flushright}
	% 	      $\to$ \textcite{Miura201705}のP.30--34，\textcite{Toyoda1992}のP.24-66.
	%       \end{flushright}

	\item[ランダム化の意義と方法] ランダム化(無作為化)は実験研究の核心部分であり，偶然誤差や交絡要因の影響を統制するために不可欠な手法である。ランダム化比較試験(RCT)では，研究参加者を無作為に実験群と統制群に割り当てることで，群間の初期状態の等質性を確保する。これにより，介入効果以外の要因による影響を最小化し，観察される群間差を介入効果として解釈できるようになる。ランダム化の具体的な方法(単純ランダム化，ブロックランダム化，層別ランダム化など)とその利点・欠点について学ぶ。
	      \begin{flushright}
		      $\to$ \textcite{Yasui2019}を参照，RCTの基本的な考え方については\textcite{Toyoda201706}のP.5-26.
	      \end{flushright}

	\item[用語の整理] 実験計画を構成する基本要素として，要因(factor)と水準(level)の概念を理解する。一要因計画と多要因計画の違い，群間計画(between-subjects design)と群内計画(within-subjects design)の特徴と選択基準について学ぶ。また，各種実験計画の表記法(例：「2$\times$3の混合計画」)についても理解する。次回以降の講義で詳細に学ぶ群間計画と群内計画の基本的な考え方について概観し，それぞれの利点と欠点を理解する。
	      \begin{flushright}
		      $\to$ \textcite{Toyoda201706}のP.27--44，実験計画の種類については\textcite{kosugi2018}のP.74--78.
	      \end{flushright}

	\item[効果はどのように評価するのか] 実験効果を評価するためには，介入による変化(効果)と偶然による変動(誤差)を区別する必要がある。効果の大きさは，群間差(実験群平均 - 統制群平均)で評価するが，その統計的有意性は誤差の大きさとの相対比較によって判断される。実験効果の評価における統計的検定の役割と，効果量(標準化された効果の指標)の重要性について理解する。
	      \begin{flushright}
		      $\to$ \textcite{kosugi2018}のP.60--67，効果量については\textcite{Matia201201}のP.43--46.
	      \end{flushright}

	\item[線形モデルとしてみた要因計画] 実験データには必ず誤差(偶然変動)が含まれる。誤差は様々な要因の集積によって生じるが，中心極限定理により，多くの独立した小さな誤差要因が加算されると，その合計は正規分布に近似する。この性質が実験データの統計的分析の理論的基盤となっている。誤差分布としての正規分布の特性(平均0，対称性，確率的性質)を理解し，実験計画における誤差の扱い方と統制方法について学ぶ。効果と誤差の分離は，線形モデルの枠組みで理解できる(個々の観測値 = 全体平均 + 介入効果 + 誤差)。平方和の分解のモデルが回帰分析と同じ形をしていること，帰無仮説を数式で表現すると水平な直線であり，対立仮説は傾きのある直線として表現できることを確認する。これにより，分散分析が回帰分析の特別なケースであることを理解する。
	      \begin{flushright}
		      $\to$ 誤差分布については\textcite{Naganuma2201611}のP.16--40，\textcite{Yamada2004}のP.80--85.
	      \end{flushright}

\end{description}
\subsubsection{キーワード}
\begin{itemize}
	\item ランダム化比較試験(RCT)
	\item 実験群と統制群
	\item 偶然誤差と正規分布
	\item 要因と水準
	\item 群間計画と群内計画
\end{itemize}
\subsection{授業情報}
\paragraph{コマの展開方法}
講義
\paragraph{公認心理師標準シラバスにおける位置づけ}
科目番号4 心理学研究法；①心理学における実証的研究法(量的研究及び質的研究)；B 実験的方法と観察的方法，D 実験的方法(1)実験法 に相当
\paragraph{心理統計法標準カリキュラム(日本心理学会)における位置づけ}
第5回 相関と因果 細目3，第6回 実験的研究と要因計画 細目1・3・4 に相当
\subsubsection{予習・復習課題}
\paragraph{予習}
実験計画法の基本的な考え方について，心理学研究法の教科書等で予習しておくこと。特に因果関係と相関関係の違い，実験研究と調査研究の違いについて理解しておくと講義の理解が深まる。また，前回学んだt検定の基本的な考え方を復習しておくと，実験計画と統計的分析の関連について理解しやすくなる。

\paragraph{復習}
授業で学んだ内容を定着させるために，以下の点について理解を深めておくこと：

\begin{enumerate}
\item ランダム化の意義と具体的な方法について説明できるようにする
\item 偶然誤差がなぜ正規分布に従うと考えられるのか，その理論的根拠を理解する
\item 実験効果と誤差を分離する考え方を理解し，線形モデルの枠組みで説明できるようにする
\item 様々な実験計画の種類とその特徴を理解し，具体的な研究例を挙げて説明できるようにする
\item 実験研究の内的妥当性と外的妥当性のバランスについて考察する
\end{enumerate}

また，次回の授業で学ぶ群間計画(Between-subjects design)の準備として，\textcite{Toyoda201706}のP.27--44や\textcite{kosugi2018}のP.74--78を読んでおくとよい。
\section{群間計画(Between Design)と分散分析}
\label{b18_BetweenDesign}

\subsection{授業内容}
\subsubsection{概要}

前回学んだ実験計画法の基礎知識を踏まえ，本コマでは群間計画(Between-subjects Design)に焦点を当てる。群間計画とは，異なる実験参加者グループに異なる条件を割り当てる実験デザインであり，心理学研究でよく用いられる手法である。特に，一要因および二要因の群間計画における効果の分解方法と評価方法を学び，分散分析(ANOVA)の基本的な考え方を理解する。平方和の分解，交互作用の概念について学習する。

\subsubsection{コマ主題細目}
\begin{description}
	\item[一要因群間計画の平方和の分解] 群間計画(Between-subjects Design)は，異なる実験参加者グループに異なる条件を割り当てる実験デザインである。群間計画では，各参加者は一条件のみを経験するため，練習効果や順序効果といった問題を回避できるが，個人差による誤差が大きくなる傾向がある。分散分析の核心は，データの総変動(平方和)を複数の変動源に分解することにある。一要因群間計画では，総平方和($SS_T$)は要因の効果による平方和($SS_A$)と誤差による平方和($SS_E$)に分解される($SS_T = SS_A + SS_E$)。各平方和の計算方法と意味を理解し，効果と誤差をどのように分離するかを学ぶ。
	      \begin{flushright}
		      $\to$ \textcite{Toyoda201706}のP.27--44，\textcite{kosugi2018}のP.74--78.
	      \end{flushright}

	\item[二要因群間計画と主効果] 二要因群間計画では，一要因の場合に加えて，第二要因($SS_B$)と交互作用($SS_{AB}$)の平方和が加わる($SS_T = SS_A + SS_B + SS_{AB} + SS_E$)。全体平均，群平均，セル平均の関係と，それらを用いた効果の算出方法についても理解する。
	      \begin{flushright}
		      $\to$ \textcite{Yamada2004}のP.184--189，\textcite{kosugi2018}のP.81--88.
	      \end{flushright}


	\item[交互作用の理解と解釈] 二要因以上の実験計画では，要因間の交互作用(interaction)が重要な意味を持つ。交互作用とは，一方の要因の効果が他方の要因の水準によって異なる現象を指す。主効果が周辺平均の比較であるのに対し，交互作用はセルの中での比較であり，いずれも全体平均からの相対的な比較であるから効果の総和はゼロになる。交互作用の種類，交互作用の平方和を算出する手順を学ぶ。交互作用の視覚的理解を助けるグラフの描き方と読み方についても習得する。
	      \begin{flushright}
		      $\to$ \textcite{Yamada2004}のP.190--193，\textcite{kosugi2019}のP.139--142.
	      \end{flushright}

	\item[平方和の分解と変動の源泉] 要因の数が増えると，平方和の分解に現れる交互作用の組み合わせが急速に増える。3要因になると，3つの主効果に加え，AB，AC，BCの一次の交互作用だけでなく，ABCの二次の交互作用も考える必要がある。要因計画はいくらでも複雑にしてよいものではなく，要因が増えるごとに考えるべき要素が増え，統一的な解釈が難しくなることを理解する。


\end{description}
\subsubsection{キーワード}
\begin{itemize}
	\item 群間計画(Between-subjects Design)
	\item 平方和の分解と分散分析
	\item 主効果と交互作用
\end{itemize}
\subsection{授業情報}
\paragraph{コマの展開方法}
講義
\paragraph{公認心理師標準シラバスにおける位置づけ}
科目番号5 心理学統計法；①心理学で用いられる統計手法；H 推測統計：代表値と散布度をめぐって(2) に相当
\paragraph{心理統計法標準カリキュラム(日本心理学会)における位置づけ}
第6回 実験的研究と要因計画 細目2，第11回 分散分析 細目1・3 に相当
\subsubsection{予習・復習課題}
\paragraph{予習}
前回学んだ実験計画法の基礎概念(要因と水準，ランダム化など)を復習しておくこと。また，線形モデルの基本的な考え方(切片と傾き，回帰式)についても再確認しておくと理解が深まる。さらに，分散(variance)の概念と計算方法について復習しておくことも有用である。

\paragraph{復習}
授業で学んだ内容を定着させるために，以下の点について理解を深めておくこと：

\begin{enumerate}
\item 一要因および二要因の群間計画における平方和の分解方法を理解し，具体例を用いて計算できるようにする
\item 交互作用の概念を理解し，様々なパターンの交互作用(相乗効果，拮抗効果，クロスオーバーなど)を図示して説明できるようにする
\item 提示された実験デザインを適切に表記し，その構造(要因，水準，群間/群内)を説明できるようにする
\end{enumerate}

次回の授業では，Rを用いた群間計画データの分析を行うため，統計環境Rの基本操作を再確認しておくことが望ましい。特に，データフレームの操作方法や基本的な関数の使い方について復習しておくと良い。
\section{F検定/分散分析の基礎}
\label{b19_Ftest}

\subsection{授業内容}
\subsubsection{概要}
前回学んだ群間計画(Between-subjects Design)の理論的基礎として，平方和の分解について理解した。これと統計的判断すなわちF検定を組み合わせることで，実際のデータに基づいて効果の有意性を評価できるようになる。Rを用いて一要因の群間計画データを実際に分析し，分散分析表の読み方や結果の解釈方法を習得する。基本的な分析には\verb|aov()|関数と\verb|lm()|関数を用い，それらの出力結果からSS(平方和)，df(自由度)，MS(平均平方)，F値などの情報をどのように読み取るかを学ぶ。また，検定の帰無仮説と対立仮説を改めて確認し，有意な結果が得られた場合に，どの群間に具体的な差があるのかを特定するための下位検定(post-hoc test)の方法についても学ぶ。

\subsubsection{コマ主題細目}
\begin{description}

	\item[検定統計量の算出] 平方和を分解した後，効果の統計的有意性を評価するためにF検定を用いる。F値は，効果の平均平方($MS = SS/df$)を誤差の平均平方で割った値(F = MS効果/MS誤差)として算出される。自由度(df)の概念と計算方法，F分布の特性とF値の解釈について学ぶ。特に，F値が大きいほど効果が誤差に比べて相対的に大きいことを意味し，これが統計的に有意であるかどうかを判断するための基準となることを理解する。
	\begin{flushright}
		$\to$ \textcite{Kawabata201408}のP.125--130，\textcite{Yamada2004}のP.190--196.
	\end{flushright}

	\item[分散分析の手順] 帰無仮説検定の手順(仮説の設定，検定統計量の選択，判断の基準になる確率の設定，検定統計量の実現値の計算，判断)を一要因の例で改めてたどる。分散分析における帰無仮説と対立仮説を確認し，何がどのように判断されたのかを改めて確認する。平方和，自由度，平均平方，F値，$p$値の流れを分散分析表(ANOVA table)にまとめて報告する方法と，効果量$\varepsilon^2$の算出方法を学ぶ。

	\item[Rで一要因群間計画の分析] 分析に先立って，Rの変数の型(numeric，factor，chrなど)を確認する。特にRでの分散分析では，カテゴリカル変数を因子型(factor)として扱うことが重要である。データフレームの作り方，\verb|str()|関数による構造の確認，因子型への変換方法を学ぶ。Rの\verb|aov()|関数と\verb|lm()|関数を用いた一要因分散分析の実行方法を学ぶ。公式表記(formula notation)の書き方(例：\verb|y ~ group|)，関数の使い方，出力結果の読み方について理解する。\verb|summary()|関数を用いた分散分析表(ANOVA table)の表示方法と，そこに含まれるSS(平方和)，df(自由度)，MS(平均平方)，F値，p値の意味を確認する。
	      \begin{flushright}
		      $\to$ \textcite{kosugi2019}のP.46--70，P.146--152，\textcite{Kinosady2021}のP.28--42，\textcite{Hashimoto201607}のP.26--55.
	      \end{flushright}

	\item[Rで下位検定]
	分散分析における下位検定(post-hoc test)の必要性と目的について学ぶ。有意な結果が得られた場合，どの群間に具体的な差があるのかを特定するために下位検定を行う必要があることを理解する。代表的な下位検定の方法(例：TukeyのHSD，Bonferroni法，Scheffé法など)について理解する。また，\verb|TukeyHSD()|関数を用いた多重比較の方法や，効果量($\eta^2$，$\omega^2$など)の算出方法についても学ぶ。実際のデータセットを用いた演習を通じて，一要因群間計画の分析手順と結果の解釈を身につける。



\end{description}
\subsubsection{キーワード}
\begin{itemize}
	\item \verb|aov()|関数と\verb|lm()|関数
	\item 分散分析表の読み方(SS, df, MS, F値)
	\item 多重比較と下位検定
\end{itemize}
\subsection{授業情報}
\paragraph{コマの展開方法}
演習
\paragraph{公認心理師標準シラバスにおける位置づけ}
科目番号5 心理学統計法；①心理学で用いられる統計手法；G 推測統計：代表値と散布度をめぐって(1)，H 推測統計：代表値と散布度をめぐって(2) に相当
\paragraph{心理統計法標準カリキュラム(日本心理学会)における位置づけ}
第11回 分散分析 細目2・5，第12回 統計的仮説検定の注意点 細目2 に相当
\subsubsection{予習・復習課題}
\paragraph{予習}
前回学んだ群間計画(Between Design)の基本概念と分散分析の理論を復習しておくこと。特に，平方和の分解，F値の意味，主効果と交互作用の概念について理解を確認しておくとよい。また，RとRStudioの基本操作，特にデータの読み込み方法やデータフレームの操作方法について再確認しておくこと。必要に応じて，第3回，第4回の授業で学んだR操作についての資料を見直しておくとよい。

\paragraph{復習}
授業で学んだ内容を定着させるために，以下の課題に取り組むこと：

\begin{enumerate}
\item 提供されたデータセットを用いて，一要因群間計画の分散分析を実行し，結果を適切に解釈する
\item \verb|TukeyHSD()|関数を用いて多重比較を行い，どの群間に有意差があるかを特定する
\end{enumerate}

これらの課題をRスクリプトにまとめ，期日までに提出すること。時間内に終わらなかった場合は，次回授業までに完成させておくこと。分析結果の適切な解釈と可視化を重視し，統計的に有意な結果だけでなく，効果量や信頼区間についても言及するよう心がけること。不明点があれば，\textcite{kosugi2019}や\textcite{Hashimoto201607}を参照すること。
\section{Rによる群間計画の分析}
\label{b20_BetweenDesign_onR}

\subsection{授業内容}
\subsubsection{概要}
前回に続いて，群間計画(Between Design)のデータ分析をRで実行する方法を学ぶ。今回は，二要因群間計画の分散分析(ANOVA)の実行方法と結果の解釈についても学ぶ。Rの\verb|aov()|関数と\verb|lm()|関数でもよいが，\verb|anovakun|関数を使うとより簡単に，詳細な結果を得ることができる。分散分析表から主効果と交互作用の検定結果を読み取り，有意な結果が得られた場合の追加分析(下位検定や単純主効果の検定)の方法についても学ぶ。また，分散分析の結果を適切に可視化する方法についても学ぶ。Rの基本グラフィックス機能やggplot2パッケージを用いた群間の比較を示す棒グラフやボックスプロット，交互作用を示す交互作用プロットの作成方法を理解する。エラーバーの追加方法，軸ラベルやタイトルの設定，凡例の調整など，読みやすく情報量の多いグラフを作成するためのテクニックを学ぶ。実際のデータセットを用いた演習を通じて，群間計画の分析手順と結果の解釈，および可視化方法を身につける。


\subsubsection{コマ主題細目}
\begin{description}
	\item[分散分析用の関数の導入]
	Rで分散分析を行う際に便利な\verb|anovakun|のインストールと基本的な使い方を学ぶ。\verb|anovakun|はCRANに登録されているようなパッケージとは違い，インターネットを経由してソースコードを環境に導入する必要がある。実験デザインの設定の方法についても理解する。

	\item[データの作りかた] 分析に先立って，データセット(ロング型・ワイド型)についての理解が必要である。データファイルは1セルに1情報，重複や集計，オブジェクトや文字列のないデータセットであることが基本であり，一行ごとに1ケースのデータが入り，各列が変数を表す整然データの考え方を学ぶ。可視化のためのロング型と，分散分析の関数に渡すワイド型の変換についても理解する。

	\item[二要因Between計画の実際] 二要因分散分析の実行方法と結果の解釈について学ぶ。\verb|aov()|関数における二要因の公式表記(例：\verb|y ~ factor1 * factor2|)の意味と使い方，交互作用項の指定方法を理解する。分散分析表から主効果と交互作用の検定結果を読み取る方法，有意な交互作用がある場合の解釈と追加分析(単純主効果の検定など)の必要性について学ぶ。また，二要因分散分析の結果を適切に報告する方法についても理解する。実際のデータセットを用いた演習を通じて，二要因群間計画の分析手順と結果の解釈を身につける。
	      \begin{flushright}
		      $\to$ \textcite{kosugi2019}のP.153--160，\textcite{Hashimoto201607}のP.57--90.
	      \end{flushright}

	\item[分析結果の可視化] 分散分析の結果を適切に可視化する方法について学ぶ。Rの基本グラフィックス機能やggplot2パッケージを用いた群間の比較を示す棒グラフやボックスプロット，交互作用を示す交互作用プロットの作成方法を理解する。エラーバーの追加方法，軸ラベルやタイトルの設定，凡例の調整など，読みやすく情報量の多いグラフを作成するためのテクニックを学ぶ。また，可視化することの重要性と，異なるグラフが伝える情報の違いについても理解する。実際のデータセットを用いた演習を通じて，分析結果の効果的な可視化方法を身につける。
	      \begin{flushright}
		      $\to$ \textcite{Kinosady2021}のP.126--138，\textcite{kosugi2019}のP.90--105.
	      \end{flushright}

\end{description}
\subsubsection{キーワード}
\begin{itemize}
	\item \verb|anovakun()|関数の導入と使い方
	\item \verb|anovakun()|関数の各種オプション設定
	\item 交互作用の可視化
	\item ggplot2による結果の表現
\end{itemize}
\subsection{授業情報}
\paragraph{コマの展開方法}
演習
\paragraph{公認心理師標準シラバスにおける位置づけ}
科目番号5 心理学統計法；①心理学で用いられる統計手法；H 推測統計：代表値と散布度をめぐって(2)／科目番号5 心理学統計法；②統計に関する基礎的な知識；E エクセル，R，SPSS入門(3) に相当
\paragraph{心理統計法標準カリキュラム(日本心理学会)における位置づけ}
第11回 分散分析 細目3・4・5，第8回 統計的推定 細目5 に相当
\subsubsection{予習・復習課題}
\paragraph{予習}
前回学んだ群間計画(Between Design)の基本概念と分散分析の理論を復習しておくこと。特に，平方和の分解，F値の意味，主効果と交互作用の概念について理解を確認しておくとよい。また，RとRStudioの基本操作，特にデータの読み込み方法やデータフレームの操作方法について再確認しておくこと。必要に応じて，第3回，第4回の授業で学んだR操作についての資料を見直しておくとよい。

\paragraph{復習}
授業で学んだ内容を定着させるために，以下の課題に取り組むこと：

\begin{enumerate}
\item 同じく提供されたデータセットを用いて，\verb|anovakun|で二要因群間計画の分散分析を実行し，主効果と交互作用の結果を解釈する
\item 分散分析の結果を適切に可視化し，群間の差や交互作用を示すグラフを作成する
\end{enumerate}

これらの課題をRスクリプトファイルにまとめ，期日までに提出すること。時間内に終わらなかった場合は，次回授業までに完成させておくこと。分析結果の適切な解釈と可視化を重視し，統計的に有意な結果だけでなく，効果量や信頼区間についても言及するよう心がけること。不明点があれば，\textcite{kosugi2019}や\textcite{Hashimoto201607}を参照すること。 
\section{群内計画(Within Design)の理論}
\label{b21_WithinDesign} 
\subsection{授業内容}
\subsubsection{概要}
これまで学んだ群間計画に続き，本コマでは群内計画(Within-subjects Design)について学ぶ。群内計画は，同一の実験参加者が複数の条件を経験する実験デザインであり，事前-事後測定や反復測定などがこれに該当する。群内計画の最大の特徴は，個人差を個別に識別し，誤差から分離できることである。このことにより分析の精度が高まり，効果の検出力が向上する。本コマでは，群内計画の基本的な考え方，データの構造，効果と誤差の分離方法，個人差の取り扱い方について理解する。

\subsubsection{コマ主題細目}
\begin{description}
	\item[群内計画の基本構造] 群内計画(Within-subjects Design)は，同一の実験参加者が複数の条件を経験する実験デザインである。このデザインの特徴，利点と欠点，適用場面について理解する。群内計画の代表的な例として，事前-事後測定，反復測定，対応のある比較などがある。個人差を誤差から分離できるという最大の利点がある一方で，練習効果，疲労効果，順序効果といった問題も生じうることを理解する。また，実験のコストや実施上の制約，倫理的配慮についても考察する。群内計画を表記する際の規則(例：「3水準の一要因Within計画」)についても学ぶ。
	      \begin{flushright}
		      $\to$ \textcite{Toyoda201706}のP.45--62，\textcite{kosugi2018}のP.74--78.
	      \end{flushright}

	\item[一要因群内モデルの平方和の分解] 群内計画のデータ構造を数学的に表現し，平方和の分解方法を理解する。総平方和($SS_T$)は，要因の効果による平方和($SS_A$)，個人差による平方和($SS_S$)，および残差による平方和($SS_E$)に分解される($SS_T = SS_A + SS_S + SS_E$)。各平方和の計算方法と意味を理解し，効果と誤差，個人差をどのように分離するかを学ぶ。個人差の分散と残差の分散を分離する方法，それによってF値がどのように変化するかを理解する。群内計画の分散分析表の読み方と，結果の解釈方法についても学ぶ。
	      \begin{flushright}
		      $\to$ \textcite{Yamada2004}のP.178--183，図7.5.1の平方和の分解図式，\textcite{kosugi2018}のP.81--88.
	      \end{flushright}

	\item[二要因群内モデルの例] 二要因の群内計画で，個人差，主効果，交互作用に加えて，要因Aに伴って生じる残差，要因Bに伴って生じる残差，交互作用に伴って生じる残差を順に算出する。多要因群内計画における効果の分解と誤差項の選択，適切なF値の算出方法について学ぶ。
	      \begin{flushright}
		      $\to$ \textcite{kosugi2018}のP.125--132.
	      \end{flushright}

	\item[群間計画と群内計画の比較] 群内計画の最大の特徴は，個人差を誤差から分離できることである。従来の群間計画では，個人差は誤差に含まれていたが，群内計画では個人要因として明示的に扱うことができる。これにより比較すべき誤差(残差)が小さくなり，効果の検出力が向上する。視覚的な例を通じて，個人差を分離することの意義を直感的に理解する。また，二要因の群内計画では，比較すべき誤差が要因ごとに「要因に伴う誤差」として分かれることを確認する。
	      \begin{flushright}
		      $\to$ \textcite{kosugi2019}のP.146--148.
	      \end{flushright}

\end{description}
\subsubsection{キーワード}
\begin{itemize}
	\item 群内計画(Within-subjects Design)
	\item 個人差の分離と平方和の分解
	\item 反復測定と対応のある比較
\end{itemize}
\subsection{授業情報}
\paragraph{コマの展開方法}
講義
\paragraph{公認心理師標準シラバスにおける位置づけ}
科目番号5 心理学統計法；①心理学で用いられる統計手法；G 推測統計：代表値と散布度をめぐって(1)，H 推測統計：代表値と散布度をめぐって(2) に相当
\paragraph{心理統計法標準カリキュラム(日本心理学会)における位置づけ}
第6回 実験的研究と要因計画 細目3，第11回 分散分析 細目1 に相当
\subsubsection{予習・復習課題}
\paragraph{予習}
前回までに学んだ群間計画(Between Design)の特徴と分散分析の基本的な考え方を復習しておくこと。特に，一要因の分散分析における平方和の分解，F値の意味について理解を確認しておくとよい。また，実験計画の用語(要因，水準など)についても再確認しておくこと。心理学の論文から群内計画を用いた研究例を探して読んでみることも有用である。

\paragraph{復習}
授業で学んだ内容を定着させるために，以下の点について理解を深めておくこと：

\begin{enumerate}
\item 群間計画と群内計画の違い，それぞれの利点と欠点を説明できるようにする
\item 群内計画における平方和の分解方法を理解し，個人差が誤差から分離される過程を説明できるようにする
\item 群内計画における適切なF値の算出方法を理解し，分散分析表を正しく解釈できるようにする
\item 実際の研究例をもとに，群内計画が適切な状況とそうでない状況を判断できるようにする
\end{enumerate}

群内計画における個人差を算出するための手続きを，具体的な数値例を用いて自分で確認してみるとよい。公開されている実験データや，教科書に載っている例題を使って練習するとよい。
\section{RによるWithinデザインの分散分析}
\label{b22_WithinDesign_onR} 
\subsection{授業内容}
\subsubsection{概要}
前回学んだ群内計画(Within-subjects Design)の理論的理解をもとに，本コマではRを用いた群内計画の実践的な分析方法について学ぶ。特に，anovakunという関数群を活用し，群内計画データの構造化，モデルの構築，分散分析の実施，結果の解釈までの一連のプロセスを習得する。群内計画特有の個人差の分離，効果の検出，交互作用の解釈についてRの出力結果から理解を深める。また，分析結果の適切な可視化と報告方法についても学び，心理統計学の実践的スキルを向上させる。

\subsubsection{コマ主題細目}
\begin{description}
	\item[一要因Within計画の実際] anovakunを用いた一要因群内計画の分析方法を習得する。群内計画を指定するための構文，データの渡し方を理解する。球状性の検定の意味と，その結果に応じた自由度の調整について学ぶ。得られた分散分析表の読み方と解釈，特に個人差(参加者間変動)が明示的に分離されていることを確認し，下位検定の結果を読み取る。効果量($\eta^2$，$\omega^2$など)の算出方法と，分析結果の適切な報告形式についても習得する。
	      \begin{flushright}
		      $\to$ \textcite{kosugi2019}のP.152--157.
	      \end{flushright}

	\item[二要因Within計画の実際] 二要因以上の群内計画の分析方法について学ぶ。複数の実験要因がすべて群内要因である場合のモデル式の記述法，交互作用の指定方法，結果の解釈方法を理解する。実際の心理学実験データを用いた演習を通じて，多要因群内計画の分析スキルを向上させる。
	      \begin{flushright}
		      $\to$ \textcite{kosugi2018}のP.125--132.
	      \end{flushright}

	\item[単純効果の分散分析表の分割] 交互作用が有意であった場合の単純主効果の検定方法についてanovakunでの実施方法を習得する。一方の要因の水準ごとに他方の要因の効果を検定し，分散分析表がどのように分割されるかを理解する。

	\item[分散分析のまとめ] 分散分析は，データのばらつきを要因ごとに分解し，その要因がデータに与える影響を検証する手法である。「分散」分析という名前がついているが，目的は平均値の比較であることを確認する。群間計画では各要因の効果と全体の誤差項を，群内計画では各要因の効果とそれに伴って生じる誤差を使って$F$値を計算することを，これまでの内容と合わせて整理する。

	% \item[混合計画の場合] 群内要因と群間要因が混在する混合計画の分析方法について学ぶ。anovakunでのモデル式の記述法，交互作用の指定方法，結果の解釈方法を理解する。特に，混合計画における個人差の分離と効果量の算出方法についても学ぶ。実際の心理学実験データを用いた演習を通じて，混合計画の分析スキルを向上させる。

\end{description}
\subsubsection{キーワード}
\begin{itemize}
	\item anovakunによる群内計画分析 
	\item 一要因および多要因群内計画の分散分析
	\item 球状性の検定
	\item 効果量の算出と解釈
\end{itemize}
\subsection{授業情報}
\paragraph{コマの展開方法}
講義とR実習
\paragraph{公認心理師標準シラバスにおける位置づけ}
科目番号5 心理学統計法；①心理学で用いられる統計手法；H 推測統計：代表値と散布度をめぐって(2) に相当
\paragraph{心理統計法標準カリキュラム(日本心理学会)における位置づけ}
第11回 分散分析 細目4，第12回 統計的仮説検定の注意点 細目3 に相当
\subsubsection{予習・復習課題}
\paragraph{予習}
前回学習した群内計画(Within Design)の理論的内容を復習しておくこと。特に，群内計画における個人差の分離の意義，平方和の分解，適切なF値の算出方法について理解を確認しておくとよい。また，Rの基本操作，データフレームの扱い方，基本的な関数の使用方法についても再確認しておくこと。

\paragraph{復習}
授業で学んだ内容を定着させるために，以下の点について理解を深めておくこと：

\begin{enumerate}
\item anovakunを用いた一要因群内計画の分析を独力で実施し，結果を正しく解釈できるようにする 
\item 多要因群内計画の分析方法を理解し，交互作用の解釈と単純効果を理解できるようにする
 \item 群内計画の分析結果を適切に可視化し，レポートで報告できる書き方を身につける
\item 実際の研究データや公開データを用いて，群内計画の分析を練習し，スキルを向上させる
\end{enumerate}

授業で扱った分析スクリプトを，異なるデータセットに適用してみることで理解を深めるとよい。公開されている心理実験データや教科書に掲載されているデータを用いて実習してみるとよい。また，各自の研究計画において群内計画が適用可能かどうかを検討し，必要に応じて分析計画を立ててみるとよい。なお，anovakunの開発者のWebサイトやGitHubページで提供されている追加の解説やチュートリアルも参考になる。
\section{確率モデルと尤度}
\label{s23_Probabilitic_model}

\subsection{授業内容}
\subsubsection{概要}
ここまでは母平均や母分散の推定に標本統計量を用いてきたが，この方法はモーメント法と呼ばれるもので，母数を推定する唯一の方法ではない。統計モデルは確率分布を用いて記述され，この確率モデルに基づいて母数を推定する方法として，最尤法とベイズ法がある。ここでは最尤法の基礎となる確率モデルの考え方と，尤度関数について理解する。
\subsubsection{コマ主題細目(箇条書き)}
\begin{description}
	\item[確率関数と尤度関数]最も単純な例として，ベルヌーイ分布を取り上げる。ベルヌーイ分布は確率$\theta$によって0または1が得られる分布であるが，推定に当たっては$\theta$が未知でデータが既知である。この関係から，関数の読み取り方を逆にしたものがとくに尤度関数と呼ばれる。

	\item[最尤推定]データに基づいて尤度を計算できることから，このデータが得られるのに最も適したパラメータの位置を推定する方法が最尤推定である。最尤推定では尤度関数の総積を計算し，そのピークの値を推定値とする方法である。尤度関数を実際に描きながら解説する。ここで注意するべき点は，尤度関数もまたなんらかの分布関数のように見えるかもしれないが，公理よりこれが確率分布関数ではないことをしっかり意識することである。
	      \begin{flushright}
		      $\to$ \textcite{Nojima201903}のP.74にベルヌーイ分布の例。あるいは\textcite{kosugi2018}のP.117--119.
	      \end{flushright}


	\item[対数尤度関数]確率密度の総積を計算すると，その値は必然的にどんどん小さなものになっていき，計算が非常に困難である。そこで実際の計算では，尤度関数の対数を取ることが一般的である。対数を取っても値の大小関係は変わらないため，ピークを推定するのには問題なく利用できる。記号が$\prod$から$\sum$に変わることは対数の特徴である。

	\item[正規分布の尤度関数]心理学では主に正規分布をモデルに用いるが，正規分布のパラメータは2つある。話を簡単にするため，まずは分散パラメータを固定し，平均パラメータを変化させればどのように尤度関数が，対数尤度関数が変化するかを確認する。つづいて，分散パラメータも未知なものとして，両者が同時にピークを迎える点を最尤推定とする。

\end{description}
\subsubsection{キーワード}
\begin{itemize}
	\item 確率密度関数，確率質量関数
	\item 尤度関数
	\item 最尤推定値
	\item 対数尤度
\end{itemize}
\subsection{授業情報}
\paragraph{コマの展開方法}
講義
\paragraph{公認心理師標準シラバスにおける位置づけ}
科目番号5 心理学統計法；②統計に関する基礎的な知識；B 確率と確率分布 に相当
\paragraph{心理統計法標準カリキュラム(日本心理学会)における位置づけ}
該当なし
\subsubsection{予習・復習課題}
\paragraph{予習}
改めて確率分布を用いたモデルに展開していくため，これまでに習った確率分布についての基礎知識(確率質量，確率密度，確率変数やパラメータを使った記述の仕方，等)について，これまでの講義を復習しておくことが望ましい。また，総和の記号$\sum$や対数についての考え方など，これまでに習っている教材を利用して復習しておくとよい。
\paragraph{復習}
非常に小さな数字の計算になるため，尤度関数を最大化する問題はコンピュータにとっても難しい。とはいえ，表計算ソフトや統計環境をもちいて，簡単な数値例を自ら考え，計算したり図示したりすることで直観的に理解が進むこともある。また対数を取っても大小関係が変わらないことなども，手元の計算機(電卓，表計算ソフト，統計環境R)で簡単に確認できるので，是非数値例を用いた復習を行っておいて欲しい。
\section{尤度と線形モデル}
\label{b24_linear_model}

\subsection{授業内容}

\subsubsection{概要}
最尤法では確率モデルを使って母数を推定するが，心理学では主に正規分布をモデルに用いる。正規分布の平均パラメータを線形予測子で表現し，誤差項が正規分布に従うと仮定することで，回帰モデルを確率モデルとして表現できることを理解する。また，回帰モデルのパラメータ(切片，傾き，誤差分散)を最尤推定する方法を学ぶ。さらに，線形モデルを一般化し，様々なタイプのデータに対応する方法として，一般化線形モデル(GLM: Generalized Linear Model)の基礎について理解する。

\subsubsection{コマ主題細目}
\begin{description}
	\item[回帰モデルの確率表現] 単回帰モデル$Y_i = \beta_0 + \beta_1 X_i + \varepsilon_i$を確率モデルとして表現し，最尤推定することを考える。正規分布の平均パラメータを線形予測子$\beta_0 + \beta_1 X_i$で表現し，誤差項$\varepsilon_i$が正規分布$N(0, \sigma^2)$に従うと仮定する。このとき，観測データ$Y_i$は平均$\beta_0 + \beta_1 X_i$，分散$\sigma^2$の正規分布に従う確率変数として表現できることを理解する。
    \item[回帰分析の最尤推定] 回帰モデルのパラメータ(切片$\beta_0$，傾き$\beta_1$，誤差分散$\sigma^2$)を最尤推定する方法を学ぶ。尤度関数の構築とその最大化により，最尤推定値を導出する手順を理解する。また，Rによる実践を通じて，最尤推定の結果と最小二乗法による推定結果が一致することを確認する。
    \item[線形モデルの拡張] 線形モデルを一般化し，様々なタイプのデータに対応する方法を学ぶ。線形結合でモデルを表現すること，不確定な誤差の記述として確率分布を用い，確率分布の平均的パラメータに線形予測子を対応させることを理解する。回帰分析はこの枠組みの一つであり，かつ，リンク関数が恒等関数である特別な場合であることを学ぶ。また，群内計画は，統計的には階層モデル(multilevel model)あるいは混合モデル(mixed model)として捉えることができる。個人差という別の分布が入った統計モデルとして理解することで，より複雑な実験デザインにも対応できる柔軟性を持つ。固定効果(fixed effect)と変量効果(random effect)の違い，個人要因を変量効果として扱う意義について学ぶ。
    \item[一般化線形モデル] 具体的な確率分布を導入し，パラメータの定義域に応じたリンク関数を用いることで，二項分布やポアソン分布などの非正規分布を扱う一般化線形モデル(GLM: Generalized Linear Model)の基礎を理解する。これにより，様々なタイプのデータに対して回帰分析を適用できることを学ぶ。
\end{description}
\subsubsection{キーワード}
\begin{itemize}
	\item 回帰モデルの確率表現
	\item 最尤推定
	\item 線形予測子
	\item リンク関数
	\item 一般化線形モデル
	\item 階層モデルと混合モデル
\end{itemize}
\subsection{授業情報}
\paragraph{コマの展開方法}
講義
\paragraph{公認心理師標準シラバスにおける位置づけ}
科目番号5 心理学統計法；①心理学で用いられる統計手法；J より高度な記述や推定を目指して に相当
\paragraph{心理統計法標準カリキュラム(日本心理学会)における位置づけ}
該当なし
\subsubsection{予習・復習課題}
\paragraph{予習}
正規分布の確率密度関数を確認し，正規分布を用いた最尤推定の基礎を復習しておくこと。特に尤度と確率の違いについて整理しておくことが肝要である。また回帰分析の基本である，誤差項が正規分布に従うという仮定の根拠，ひいてはデータの背後にある確率モデルの考え方についても再確認しておくとよい。
\paragraph{復習}
  線形予測子とリンク関数の役割を整理し，一般化線形モデルがどのようにして様々なデータ型に対応できるか，自分で説明できるようになることが望ましい。また，最尤推定の手順を再確認し，Rを用いて回帰モデルの最尤推定を実践できるようにしておくこと。最小二乗法と最尤法の関係についても理解を深めておくとよい。
\section{もうひとつの推定: ベイズの定理}
\label{b25_bayesian_intro}
\subsection{授業内容}
\subsubsection{概要}
母数を推定する第三の方法として，ベイズ法による推定を考える。ベイズの定理の歴史は古く，逆確率の法則とも呼ばれるものである。ベイズの定理を理解するにあたっては，条件つき確率について理解する必要があり，これらを通じてベイズの定理を導く。ベイズ法の利点は，尤度関数を確率分布関数に変えてくれることであり，推定値の確からしさを確率で表現することができる点である。これによって前期から学んできたモーメント法・最尤法と合わせて3つの推定方法が出そろうことになるので，それぞれの前提と得意分野を一望して比較する。最後に病気の罹患率の例を用いて，確率を信念の更新として理解することもできることを学ぶ。
\subsubsection{コマ主題細目(箇条書き)}
\begin{description}
	\item[ベイズの定理の導出]これまで扱ってきた確率分布は単一のものであったが，二変数による(二次元的な)確率分布を導入する。複数の事象が同時に生じることを表す同時確率，同時確率のうち一方の変数だけに周辺化して考える周辺確率，一方の情報が分かっている上での確率である条件つき確率，の3つの考え方が導かれる。また，両者が独立であるとはどういうことかについて確率的な定義を与える。これらの用語を踏まえ，条件付き確率の式変形を通じてベイズの定理を導出する。簡単な数値例を用いて解説する。
	      \begin{flushright}
		      $\to$ \textcite{Maeda2017}のP.93--109.\textcite{Toyoda201506}のP.3--11.
	      \end{flushright}

	\item[ベイズの定理の意味] ベイズの定理について，数式をさまざまな角度から眺めてその意味を理解する。同時確率，周辺確率の定義などを用いて，条件付き確率の式を変形していくことができ，最終的に右辺と左辺で条件が逆になる場合の確率が計算できることがわかる。これがベイズの定理と呼ばれるものであり，逆確率の定理，法則と呼ばれることもある。ベイズの定理は，事前確率に尤度と周辺尤度の比をかけたものとして理解することもできるし，尤度関数を確率として捉えるために変形させる項として事前確率と周辺尤度の比をかけたもの，と考えることもできる。
	      \begin{flushright}
		      $\to$ \textcite{Maeda2017}のP.93--109.尤度を確率分布に変えることの分かりやすい説明としては，\textcite{Lambert201805}のTable4.1.Xが良い。
	      \end{flushright}

	\item[3つの推定方法の比較]ベイズ法が導入されたことで，前期から学んできたモーメント法・最尤法とあわせて3つの推定方法が出そろったことになる。ここで一度立ち止まり，それぞれを一望して比較する。モーメント法は標本統計量をそのまま母数の推定値として用いる古典的な方法で，手続きがパターン化しやすく帰無仮説検定と相性が良かったために広く普及した。最尤法は確率分布をデータ生成器と捉え，手元のデータが最も生じやすいパラメータを探す方法で，確率モデルの王道といえる位置にあるが，モデルが複雑になると計算が困難になる。ベイズ法は原理そのものは18世紀に遡るが，計算機科学の発展によって実用可能となり，事前分布という「答えを出すためのヒント」を持つことで，少ないデータや複雑なモデルでも推定を可能にする。それぞれの歴史的背景，前提，得意分野，限界を整理し，この章の章題である「もうひとつの推定」の意味を確認する。
	      \begin{flushright}
		      $\to$ \textcite{Lee201709}のP.3--13.
	      \end{flushright}

	\item[ベイズの定理を使った推論の例]ベイズの定理について，意味的な解釈を膨らませておこう。ベイズの定理は確率のアップデートとして考えることができる。我々が推論するというのは，確率をアップデートすることだと考えることもできる，ということについてはすでに触れており，確率という数字が従うべきルールさえ守られていれば，確率は人間の自然な推論モデルとして考えることができるだろう。ただし頻度主義的な確率の理解と，主観確率による理解の違いには改めて注意が必要である。ここではベイズの定理を使っての推論の例として，病気の罹患率の例を取り上げる。ポイントは，「事前の確率がわかっているということ」と「情報がアップデートされる」ということを意識するかどうかにある。なお同種の例題としてよく知られるモンティ・ホール問題は，本論の流れを妨げないよう章末の補遺に収めてあるので，各自で読み解いてほしい。
	      \begin{flushright}
		      $\to$ \textcite{Maeda2017}のP.15--23.あるいは\textcite{Toyoda201506}のP.8--12.
	      \end{flushright}

\end{description}
\subsubsection{キーワード}
\begin{itemize}
	\item 同時確率
	\item 周辺確率
	\item 条件付き確率
	\item ベイズの定理
	\item 事前確率，事後確率
	\item モーメント法・最尤法・ベイズ法の比較
\end{itemize}
\subsection{授業情報}
\paragraph{コマの展開方法}
講義
\paragraph{公認心理師標準シラバスにおける位置づけ}
科目番号5 心理学統計法；①心理学で用いられる統計手法；J より高度な記述や推定を目指して に相当
\paragraph{心理統計法標準カリキュラム(日本心理学会)における位置づけ}
該当なし
\subsubsection{予習・復習課題}
\paragraph{予習}
ここからは，より進んだ統計的内容に入る。ベイズ統計学を心理学領域で活用する歴史はまだ浅いが，現代社会では実践的に多用されている技術であり，心理学にとっても今後ますます重要性が高まることは間違いない。たとえば\textcite{Nojima201903}は公認心理師向けのテキストであるが，その最終章をベイズにあてている。
内容的には，確率についてまた違った視点からの考え方が導入される。とはいえ，確率が守るべきルールについては変わりがない。改めて第\ref{b10_Probability}講の確率について復習しておくとよい。あわせて，前期に学んだモーメント法と最尤法の考え方を思い出しておくと，本コマ後半の比較が理解しやすい。
「ベイズ統計学」をキーワードとしたいくつかの入門書があるので，目を通しておくと良い。あるいは章末の補遺に収めたモンティ・ホール問題に目を通しておくと，興味を覚える人もいるかもしれない。
\paragraph{復習}
さまざまな確率の書き方から，ベイズの定理の導出まで，筋道をもう一度自らフォローアップして導出できるかどうか確認しておくこと。罹患率のような具体例については，読み物として\textcite{Gerd200309}があるので目を通すと良い。
また，3つの推定方法(モーメント法・最尤法・ベイズ法)について，それぞれの考え方の違いと得意・不得意を自分の言葉で説明できるようにしておくこと。
確率についての説明の仕方については，\textcite{Hiraoka200910}のP.1--69の説明も異なる切り口からでおもしろいし，\textcite{Ishida201801}のP.1--12もわかりやすい。ベイズ統計学についてのごく簡単な入門書としては，ラノベ風の\textcite{Ishida201902}なども良いだろう。

\section{分布で推定する}
\label{b26_bayesian_approach}

\subsection{授業内容}
\subsubsection{概要}
ベイズの定理が導入され，確率に対するもう一つの考え方を導入した。
ここでは具体的な例を使って，事前分布で確信度を表現すること，事後分布も確率分布であり，頻度主義的な解釈とは異なることを確認する。
事後分布は確率分布なので，推定にあたっては分布を要約する必要があり，EAP, MAP, MED推定値や確信区間，HDIなど分布を要約して表現する方法について整理する。
最後に，事後分布が事前分布と尤度の折衷として形成されることを確認し，事前分布の置き方が結果に与える影響について整理する。
\subsubsection{コマ主題細目(箇条書き)}
\begin{description}
	\item[事前分布と事後分布] ここまで事前・事後確率と呼んでいたものは，連続的な分布関数であっても同様に機能する。ここで確率のベイズ的解釈により，事前分布はデータが得られる前の確信度を表現していること，またとくになんの情報もないとき，わからないことを一様分布のように表現できることを理解する。また，ベイズの定理は尤度を(事後)分布に変換するものとして捉え，結果が分布として得られることを確認する。具体例としてベルヌーイ分布をもちいた学力の推定をとりあげ，最尤推定との違いや，分布による結果であることの意味を理解する。
	      \begin{flushright}
		      $\to$ \textcite{Maeda2017}のPp.18-23.や，\textcite{Toyoda201606}のP.18--21
	      \end{flushright}

	\item[事後分布を代表する数字]ベイズによる推定の結果は，事後分布という確率分布の形で得られる。結果が分布として得られるため，報告する場合はたいていその代表値を使って報告することになる。とくにMCMCによるベイズ推定の場合は，多くのMCMCサンプルの標本記述統計量がその代表値として考えられるという利点がある。中央値(MED)，平均値(EAP)，最頻値(MAP)の名称と算出方法について理解し，そのほかにも最高密度区間(HDI)といった区間で報告することもある。
	      \begin{flushright}
		      $\to$ \textcite{Toyoda201506}のP.51--55.,HDIについては\textcite{Maeda2017}のP.91--93.
	      \end{flushright}

	\item[事前分布と尤度の折衷]事前分布を一様分布においた場合，事後分布の形は尤度関数そのままの形で現れ，かつ確率分布として考えられる利点がある。しかし事前分布が異なれば，事後分布の形はそれに影響を受け，結果が変わってしまう。事後分布は事前分布と尤度関数との折衷であり，多くのデータをもつ鋭い尤度関数であればその影響は小さいが，緩やかな尤度関数であれば大きく影響を受けてしまう。研究に際して適切な事前分布を置くことの重要性を確認する。
	      \begin{flushright}
		      $\to$\textcite{Maeda2017}のP.136--138
	      \end{flushright}

\end{description}
\subsubsection{キーワード}
\begin{itemize}
	\item 事前分布，事後分布
	\item MED, EAP, MAP, HDI
	\item 事前分布と尤度の折衷
\end{itemize}
\subsection{授業情報}
\paragraph{コマの展開方法}
講義
\paragraph{公認心理師標準シラバスにおける位置づけ}
科目番号5 心理学統計法；①心理学で用いられる統計手法；J より高度な記述や推定を目指して に相当
\paragraph{心理統計法標準カリキュラム(日本心理学会)における位置づけ}
該当なし
\subsubsection{予習・復習課題}
\paragraph{予習}確率の考え方やベイズの定理についての数式的理解，意味的理解を踏まえ，これを心理統計に応用する段階に進むところであるから，これらのについて不明な点がないように，事前に十分理解しておく必要がある。
\paragraph{復習}事後分布を要約する4つの方法(MED, EAP, MAP, HDI)について，それぞれが分布のどこを見ているのかを説明できるようにしておくこと。また，事後分布が事前分布と尤度の折衷であることを踏まえ，事前分布の置き方が結果にどう影響するのか，データ数との関係とあわせて整理しておくことが重要である。頻度主義的な考え方との違いに留意しつつ，結果が「分布として得られる」ことの意味を自分の言葉で説明できるようにしておくとよい。

\section{ベイズ的判断の方法}
\label{b27_bayes_judgement}

\subsection{授業内容}

\subsubsection{概要}
ベイズ的な判断方法は，従来の頻度主義的な$p$値に基づく仮説検定とは異なる視点からデータを解釈し，研究者がより柔軟かつ情報量豊富な判断を下すための手段を提供する。本コマでは，まず帰無仮説検定のロジックを振り返り，ベイズ的判断との対比を確認したうえで，ベイズ推定の結果をもとにした具体的な判断方法を学ぶ。すなわち，確信区間とROPEに基づく判断，そしてベイズファクターを用いた仮説検証である。最後に，これらの判断方法がもたらす利点を帰無仮説検定との対比で整理し，効果をそのままの大きさで語れること，帰無仮説を積極的に支持できること，多重比較のような補正を必要としないことを確認する。
\subsubsection{コマ主題細目}
\begin{description}
	\item[ベイズ法と帰無仮説検定] あらためて帰無仮説検定のロジックを振り返りつつ，ベイズ推定の考え方との相違について確認する。$p$値が「帰無仮説のもとでデータが得られる確率」であって「帰無仮説が正しい確率」ではないという基本的な誤解を整理し，ベイズ的アプローチでは「パラメータがどこにあるか」を直接に確率で語れるという根本的な違いを理解する。特に事前分布を活用した推定方法は，時として何の前提も置かない帰無仮説検定よりも有用であることに言及する。

	\item[確信区間に基づく判断] ベイズ推定では，パラメータの事後分布から確信区間（Credible Interval）を算出できることを学ぶ。例えば95\%HDIは，「パラメータが区間内に95\%の確率で存在する」と直接的に解釈できる。これは頻度主義的な信頼区間とは解釈が異なり，より直感的である。また，実質的等価性の領域（ROPE: Region of Practical Equivalence）の概念を導入し，パラメータがゼロや特定の値と「実質的に等価」かどうかを判断する方法を学ぶ。具体的には，95\%HDIがROPEに完全に含まれる場合は「実質的に等価」，完全に外れる場合は「実質的に異なる」，部分的に重なる場合は「判断保留」と判断できる。これにより，単なる「有意/非有意」の二分法ではなく，パラメータの大きさと不確実性を考慮した判断が可能になることを理解する。 

	\item[ベイズファクター] ベイズ統計学における仮説検証の方法として，ベイズファクター（Bayes factor, BF）の概念と解釈について学ぶ。まずベイズの公式の分母である周辺尤度$p(D \mid \mathcal{M})$が，モデル$\mathcal{M}$がどれほどデータに支持されているかを表す量であることを，パラメータの周辺化を通じて理解する。ベイズファクターはこれを二つのモデルで比べた比率であり，$BF_{01} = \frac{p(D \mid \mathcal{M}_0)}{p(D \mid \mathcal{M}_1)}$と定義される。$BF_{01} > 1$のときモデル1が，$BF_{01} < 1$のときモデル2が相対的に強く支持される。ここでモデル1を対立仮説，モデル2を帰無仮説とすれば，帰無仮説を積極的に支持することもできる。これは帰無仮説検定の枠組みでは決してできなかったことである。一方で，$p$値の5\%と同じように特定の値を目標にすることの危うさ，事前分布の置き方によって値が大きく変動すること，複雑なモデルでは周辺尤度の計算が困難になることにも注意を促す。実際の計算と実践は次コマのJASPに譲る。
	      \begin{flushright}
		      $\to$ \textcite{Lee201709}のP.88--103.
	      \end{flushright}

	\item[帰無仮説検定に対するベイズ法の利点] 以上の判断方法を通じて得られる，ベイズ推定の実践的な利点を帰無仮説検定との対比で整理する。第一に，効果をそのままの大きさで語れることである。$p$値のように有意か否かの二値に潰さず，差や効果の大きさを事後分布として直接に報告できる。第二に，帰無仮説を積極的に支持できることである。頻度主義的検定では「棄却できなかった」としか言えないが，ROPEやベイズファクターを用いれば「差がないことの証拠がどれだけあるか」を示すことができる。第三に，多重比較のような補正を必要としないことである。頻度主義的検定では検定の反復にともなって第一種の過誤が累積するため補正が要求されるが，ベイズ推定はそもそも「何回検定したか」に依存しない。これらは単なる計算手法の違いではなく，研究の報告の仕方そのものを変える視点であることを理解する。

\end{description}
\subsubsection{キーワード}
\begin{itemize}
	\item 帰無仮説検定との比較
	\item 確信区間とROPE
	\item ベイズファクター
	\item 帰無仮説検定に対するベイズ法の利点(効果量・帰無仮説の支持・補正不要)
\end{itemize}
\subsection{授業情報}
\paragraph{コマの展開方法}
講義
\paragraph{公認心理師標準シラバスにおける位置づけ}
科目番号5 心理学統計法；①心理学で用いられる統計手法；J より高度な記述や推定を目指して に相当
\paragraph{心理統計法標準カリキュラム(日本心理学会)における位置づけ}
該当なし
\subsubsection{予習・復習課題}
\paragraph{予習}
前回までに学んだベイズ推定の基礎（事後分布，確信区間，HDIなど）を復習しておくこと。特に以下の点について確認しておくとよい。
\begin{itemize}
  \item 事前分布と事後分布の関係，およびベイズの定理による更新の仕組み
  \item 事後分布の要約方法（EAP, MAP, MED, HDI）の意味と解釈
  \item 頻度主義的な信頼区間とベイズ的な確信区間の解釈の違い
\end{itemize}
また，帰無仮説検定における$p$値の解釈の問題点（例：$p$値は「帰無仮説が正しい確率」ではない）について振り返っておくと，ベイズ的判断の利点がより理解しやすくなる。参考文献として\textcite{Lee201709}のP.88--103を事前に目を通しておくことが望ましい。

\paragraph{復習}
授業で学んだ内容を定着させるために，以下の点について理解を深めておくこと：
\begin{enumerate}
  \item 確信区間とROPEに基づく判断方法について，「実質的に等価」「実質的に異なる」「判断保留」の3つの判断がどのような場合に下されるか，具体例を挙げて説明できるようにする
  \item ベイズファクター$BF_{01}$の定義と解釈について，$BF_{01}=3$や$BF_{01}=0.1$の場合にそれぞれ何を意味するか説明できるようにする 
  \item ベイズファクターに特定の目標値を設けることが，$p$値の5\%と同じ問題を招きうることを説明できるようにする
  \item ベイズ推定の3つの利点（そのままの大きさで語れる／帰無仮説を支持できる／補正がいらない）について，頻度主義的検定では何が問題になるのかとあわせて説明できるようにする
\end{enumerate}
次回のJASPを用いた実習に備えて，JASPを各自のPCにインストールしておくこと（https://jasp-stats.org/）。

\section{JASPでベイズ分析を体験する}
\label{b28_bayesian_with_jasp}

\subsection{授業内容}

\subsubsection{概要}
前コマまでにベイズ的な判断の方法を学んだが，これを実際に行うにはソフトウェアの助けが要る。プログラミングの習得を前提とせずにベイズ的アプローチを試せる道具として，本コマではJASPを用いる。JASPはオープンソースの統計ソフトウェアで，心理学研究に特化した機能と直感的なユーザーインターフェイスを提供し，頻度主義的方法とベイズ的方法を同じメニューから選べる点に特徴がある。従来通りの分析をしつつ，ベイズでやると何がどう変わるのかを比較しながら進められるわけである。本コマでは，JASPの基本的な操作方法から始め，記述統計と可視化，回帰分析におけるベイズ推定，そしてベイズファクターを用いた検定について実践的に学ぶ。これにより，学生はベイズ統計学の理論と実践の両面から理解を深め，心理学研究におけるデータ分析能力を向上させることが期待される。

\subsubsection{コマ主題細目}
\begin{description}
	\item[JASPの導入とインターフェイス] JASPはオープンソースの統計ソフトウェアで，心理学研究のための統計分析機能と直感的なユーザーインターフェイスを提供する。まず，JASPのウェブサイト（https://jasp-stats.org/）からソフトウェアをダウンロードし，各自のPCにインストールする方法を学ぶ。続いて，JASPの基本的なインターフェイスと機能（メニュー構成，データビュー，分析ビュー，結果ビュー）について理解する。JASPの特徴として，オープンサイエンスをサポートする機能（分析結果の再現性，OSFとの連携）や，従来の頻度主義的分析とベイズ的分析を同じメニューから実行できる利便性について学ぶ。また，サンプルデータの読み込み方，外部データ（CSV，SPSSなど）のインポート方法，変数の種類の設定方法についても実習を通じて習得する。
	      \begin{flushright}
		      $\to$ \textcite{YUnaShimizu2022}のP.1--15.
	      \end{flushright}

	\item[記述統計と可視化] JASPを用いた基本的なデータ分析として，記述統計と可視化の方法を学ぶ。「Descriptives」メニューを使って，各変数の基本統計量（平均，中央値，標準偏差，四分位数など）を計算し，結果の読み方を理解する。また，ヒストグラム，箱ひげ図などのグラフを作成し，データの分布特性を視覚的に把握する方法を学ぶ。これらの基本操作を通じて，JASPのワークフローに慣れ，次のステップであるより高度な分析への準備を整える。
		\begin{flushright}
		      $\to$ \textcite{YUnaShimizu2022}のP.1--15.
	      \end{flushright}

	\item[JASPによる相関・回帰分析] JASPを用いた回帰分析を，頻度主義的アプローチとベイズ的アプローチの両方で実施する。「回帰」メニューの「線形回帰」を伝統的手法とベイズ的手法，それぞれ使用して分析を行い，結果の違いを比較する。ベイズ回帰分析では，回帰係数の事後分布，確信区間，モデル比較のためのベイズファクターが出力される。

	\item[ベイズファクターを使った検定] JASPを使った平均値差の検定は，ベイズファクターを用いたモデル比較の形で行われる。例えば，2群の平均値差を検定する場合，帰無仮説モデル（平均値が等しい）と対立仮説モデル（平均値が異なる）を設定し，ベイズファクター$BF_{10}$を計算する。$BF_{10}$の値に基づいて，どちらのモデルがデータによりよく適合しているかを判断する。JASPでは，この分析を簡単に実行でき，結果の解釈方法も学ぶ。

\end{description}
\subsubsection{キーワード}
\begin{itemize}
	\item JASP
	\item GUIによるベイズ分析
	\item ベイズ回帰分析
	\item ベイズファクターによる検定
\end{itemize}
\subsection{授業情報}
\paragraph{コマの展開方法}
講義
\paragraph{公認心理師標準シラバスにおける位置づけ}
科目番号5 心理学統計法；①心理学で用いられる統計手法；J より高度な記述や推定を目指して に相当
\paragraph{心理統計法標準カリキュラム(日本心理学会)における位置づけ}
該当なし
\subsubsection{予習・復習課題}
\paragraph{予習}
前回までに学んだベイズ的判断の方法（確信区間，ROPE，ベイズファクター）について復習しておくこと。特に以下の点を確認しておくとよい：
\begin{itemize}
  \item ベイズファクター$BF_{10}$の定義と解釈基準（1--3: わずかな証拠，3--10: 中程度，10以上: 強力な証拠）
  \item 確信区間とROPEに基づく「実質的等価性」の判断方法
  \item 頻度主義的な回帰分析の基本（回帰係数，決定係数$R^2$，$t$検定など）
\end{itemize}
また，JASPを各自のPCにインストールしておくこと（https://jasp-stats.org/）。インストールが完了したら，起動して基本的なインターフェイスを確認しておくと，授業での実習がスムーズに進む。参考文献として\textcite{YUnaShimizu2022}のP.1--15を事前に読んでおくことが望ましい。

\paragraph{復習}
授業で学んだ内容を定着させるために，以下の点について理解を深めておくこと：
\begin{enumerate}
  \item JASPの基本操作（データの読み込み，変数の設定，分析メニューの選択，結果の読み方）を実際に操作して確認する
  \item 回帰モデルのベイズ表現について，パラメータ$\beta_0$，$\beta_1$，$\sigma^2$に事前分布を設定する意味を説明できるようにする
  \item 頻度主義的回帰分析とベイズ回帰分析の結果の違い（点推定値，区間推定，仮説検証の方法）を比較し，それぞれの特徴を理解する
  \item JASPで出力されるベイズファクターの読み方と解釈方法を，具体的な数値例を用いて説明できるようにする
  \item サンプルデータを用いて，JASPでの記述統計，可視化，回帰分析を一通り実行し，操作に習熟しておく
\end{enumerate}
次回はGUIでは手の届かない領域に進む。JASPのメニューにない分析をやりたくなったときにどうするか，という問題意識を持っておくと，次コマの内容が理解しやすい。

\section{ベイズモデリングへの入り口: Stanとそのフロントエンド}
\label{b29_bayesianModeling}

\subsection{授業内容}

\subsubsection{概要}
ベイズ的判断の方法を学んだあと，JASPでベイズ分析を体験した。しかしJASPのメニューにないモデルを扱いたくなったときには，確率的プログラミング言語による自前のモデリングが必要になる。本コマではまず，ベイズ法が長く実用にならなかった歴史とMCMC法の登場を振り返り，それを実装した確率的プログラミング言語がベイズ推定を実用のものにしたことを理解する。続いてその代表であるStanと，そのRフロントエンドであるbrmsを紹介し，従来の式記法から自動でStanコードが生成されることを学ぶ。さらにStan言語そのものの文法とブロック構造を学び，自分でモデルを書くための第一歩を踏み出す。最後に，確率的プログラミング言語のもとで初めて自由に設計できる生成量と事後予測分布を扱い，仮説の立て方そのものが広がることを確認する。本コマで扱うのは入口部分にとどめ，各種モデリングの具体的な応用や，心理学領域での発展的なモデリングについては3年次「ベイズ統計学」に委ねる。

\subsubsection{コマ主題細目}
\begin{description}
	\item[ベイズを実践する方法] ベイズの定理は18世紀に遡るにもかかわらず，長らく実用的な道具ではなかった。その理由は，事後分布を解析的に求めることが一般には困難で，共役事前分布を探すか数値的に総当たりするしかなかったからである。この状況を変えたのが，乱数によって事後分布からの標本を得るマルコフ連鎖モンテカルロ法(MCMC)である。ここではMCMC法の基本的な仕組みを直観的に理解する。これまでベイズ推定の章で「裏で動いている計算」と呼んできたものの正体がここで明らかになる。そのうえで，MCMC法を実装した確率的プログラミング言語(Probabilistic Programming Language; PPL)の登場によって，データと尤度と事前分布の関係を記述するだけで事後分布が得られるようになったことを学ぶ。JAGS，BUGS，Stanといった代表的な言語を紹介し，Stan言語とR言語の連携の形を確認する。

	\item[brmsによる実用的モデリング] Stanを直接書くのは敷居が高いが，Rのbrmsパッケージを使えば従来のlm()やglm()のような式記法でStanのモデルを書ける。R式から自動でStanコードとRデータが生成され，MCMCが走り，事後分布が返ってくる。具体例として，これまでの講で最小二乗法および最尤法によって推定してきた回帰分析を，同じデータのままbrmsでベイズ推定する。lm()をbrm()に書き換えるだけでよいこと，そして得られる点推定値が最尤推定とほとんど変わらないことを確認したうえで，出力の読み方の違いに注目する。標準誤差や$p$値の代わりに，事後分布の標準偏差と確信区間が並ぶのである。事後分布が乱数の集まりとして手元にあるため，「傾きがある値より大きい確率はどれほどか」といった問いに直接答えられることを理解する。

	\item[stan言語を使ったプログラミング] Stan言語はC++の流儀を引き継いだ言語であり，Rとは異なる三つの約束事がある。行末をセミコロンで区切ること，変数の型を宣言してから使うこと，そしてブロック構造で記述することである。型宣言は面倒に見えるが，標準偏差パラメータに負の値が入らないよう制約できるなど，安全装置としても働く。そのうえで，Stanのコードがdata・parameters・modelを中心とする六つのブロックから構成されることを学ぶ。題材にはRStudioが用意しているStanファイルのひな型を用い，正規分布モデルのコードを読み解く。modelブロックの記述が確率モデルの数式ほぼそのままであり，尤度関数の総乗記号を書き下す必要がないことを確認する。あわせて，事前分布を省略するとStanが無情報事前分布を仮定すること，尤度は本来データの要素ごとの積であって説明変数によって平均が変わる場合はfor文で書く必要があることを補足する。

	\item[生成量による新しい仮説の立て方] 確率的プログラミング言語のもとでは，事後分布から好きな量を計算する「生成量」を自由に設計できる。MCMCが返すのはパラメータの標本であるから，それらを組み合わせて新しい量を作れば，その事後分布がそのまま得られるのである。例えば2群の平均値の差そのもの，ある群が別の群より大きくなる確率(優越率)，両群がどの程度重なるか(被覆率)などがそれにあたる。また推定されたパラメータから新しいデータを生成すれば事後予測分布が得られ，母数ではなく「次に得られるデータ」についての判断ができるようになる。これらの指標によって，帰無仮説検定では扱えなかった「次に出会う一人」についての判断が可能になり，仮説の立て方そのものが広がることを理解する。各種モデリング(分散の検定，打ち切り・切断，混合分布，潜在クラス，階層モデルなど)の発展的な内容や心理学領域での応用例は，3年次「ベイズ統計学」(半期)で扱う。本コマはそこへの橋渡しと位置づける。
\end{description}

\subsubsection{キーワード}
\begin{itemize}
	\item MCMC法
	\item 確率的プログラミング言語
	\item Stan, brms
	\item Stanの文法とブロック構造
	\item 生成量と事後予測分布
	\item 優越率と被覆率
\end{itemize}
\subsection{授業情報}
\paragraph{コマの展開方法}
講義と事例紹介
\paragraph{公認心理師標準シラバスにおける位置づけ}
科目番号5 心理学統計法；①心理学で用いられる統計手法；J より高度な記述や推定を目指して に相当
\paragraph{心理統計法標準カリキュラム(日本心理学会)における位置づけ}
該当なし
\subsubsection{予習・復習課題}
\paragraph{予習}
JASPまでで学んだベイズ的判断の方法と，JASPの限界（メニューにないモデルは書けない）について確認しておくこと。Rとbrmsを各自の環境でインストールしておくと，授業中のデモがそのまま手元で動かせる。brmsの導入にはStanの実行環境が必要になるので，章末の補遺にあるOSごとの手順を参照し，時間に余裕をもって準備しておくこと。

\paragraph{復習}
授業で学んだ内容を定着させるために，以下の点について理解を深めておくこと：
\begin{enumerate}
  \item MCMC法が何を解決したのか，なぜそれ以前はベイズ法が実用的でなかったのかを説明できるようにする
  \item brmsがR式からStanコードを自動生成し，コンパイルしてMCMCを走らせていることを説明できるようにする
  \item 自分が知っている統計分析（t検定，回帰，分散分析）をbrmsで書いてみる
  \item Stanコードのdataブロック・parametersブロック・modelブロックがそれぞれ何を宣言する場所なのかを説明できるようにする
\end{enumerate}
本コマで内容的な授業は終了である。ベイズモデリングの世界は半期ではとても語り尽くせない。3年次の授業に進むかどうかにかかわらず，この入口で得た見取り図を元に，自分のデータと向き合う際にベイズ的アプローチも選択肢に入れられるようにしておいてほしい。

\section{後期のまとめと試験}
\label{b30_final_exam}

\subsection{授業内容}

\subsubsection{概要}
これまでの授業を通じて，心理統計学の基礎から確率モデル，そしてベイズ統計学までを学んできた。本コマでは後期の内容を総括するとともに，学習内容の理解度を確認するための授業内試験を実施する。試験は前期から後期にわたる授業で学んだ基本的な知識，重要な統計概念，データ分析の方法論，そして心理統計学の考え方の筋道について総合的に問うものであり，一年間の学習の成果を測定する。試験の内容は，記述統計から推測統計，回帰分析，分散分析，そして近年重要性を増しているベイズ統計学までを含む幅広いものとなる。

\subsubsection{コマ主題細目}
\begin{description}
	\item[後期の総括] 後期は帰無仮説検定の基礎から始まり，$t$検定，分散分析と要因計画を学んだ。ここまでは，得られたデータを確率分布に当てはめて考えるものであった。授業はそのあと，より積極的に確率分布を使う確率モデルの世界に入り，尤度と最尤法，線形モデルを経て，第三の推定法であるベイズ法に至った。事前分布と事後分布，ベイズ的な判断の方法，JASPによる実践，そして確率的プログラミング言語によるモデリングまでを一望し，「確率モデルでものを考える」という一本の筋が通っていたことを確認する。
	      \begin{flushright}
		      $\to$ これまでの各回の内容を復習しておくこと。
	      \end{flushright}

	\item[学習到達度の確認試験] 後期で学んだ内容の理解度を確認するため，授業時間内に試験を実施する。試験はマークシート形式で行われ，出題範囲は前期から後期までの全授業内容を含む。特に(1)記述統計と確率の基礎概念，(2)推測統計と仮説検定の論理，(3)分散分析と実験計画法，(4)回帰分析と相関分析，(5)一般線形モデル，(6)ベイズ統計学の基礎，(7)統計ソフトウェア（RとJASP）の操作と結果の解釈，といった領域から出題される。問題形式としては，基本的な用語や概念の理解を問う問題，計算問題，結果の解釈を問う問題，適切な分析手法の選択を問う問題などが含まれる。ソフトウェアの操作そのものではなく，分析結果の読み方や解釈，適切な分析手法の選択についての理解を問う。単なる暗記ではなく，統計的思考力を測定することを目的としている。
	      \begin{flushright}
		      $\to$ 試験の詳細については，別途配布される「試験に関するお知らせ」を参照のこと。
	      \end{flushright}

	\item[これからの学びに向けて] 心理統計の世界は広大であり，半年や一年で全てを理解できるわけではない。重要な概念として，(1)記述統計と推測統計の区別，(2)確率分布と標本分布の関係，(3)仮説検定の論理と$p$値の正しい解釈，(4)効果量と検定力の概念，(5)モデリングの考え方と一般線形モデルの枠組み，(6)頻度主義統計学とベイズ統計学の違いなどは，卒業研究や大学院での研究活動においても基盤となる。わからないことがあれば，わかるところまで戻って知識を積み上げればよい。心理統計は確実に積み上げることができる領域であり，入り口の恐怖心さえ克服すれば着実に前進できることを確認して締めくくる。
	      \begin{flushright}
		      $\to$ 3年次「ベイズ統計学」など，発展的な科目への接続を意識しておくこと。
	      \end{flushright}

\end{description}
\subsubsection{キーワード}
\begin{itemize}
	\item 心理統計学の総括
	\item 確率モデルという視点
	\item 統計的思考力
	\item 分析結果の解釈
	\item 頻度主義統計学とベイズ統計学
\end{itemize}
\subsection{授業情報}
\paragraph{コマの展開方法}
後期内容の総括(約15分)の後，授業時間内試験(約60分)を実施する
\paragraph{公認心理師標準シラバスにおける位置づけ}
科目番号5 心理学統計法の総括
\paragraph{心理統計法標準カリキュラム(日本心理学会)における位置づけ}
総括
\subsubsection{予習・復習課題}
\paragraph{予習}
本コマは総括と授業内試験であるため，以下の内容について十分な準備をして臨むこと：

\begin{enumerate}
\item 前期および後期の授業内容全体を体系的に復習し，重要な概念と用語を理解する
\item 配布資料や教科書の内容を再確認し，特に重要な統計手法とその適用条件について整理する
\item これまでに行った実習や課題の内容を振り返り，分析結果の解釈方法を確認する
\item 試験の詳細（持ち込み可能な資料の範囲など）については，別途配布されるお知らせを参照すること
\item 過去の小テストや課題で理解が不十分だった部分を特に重点的に復習する
\end{enumerate}

\paragraph{復習}
試験終了後は，問題内容を振り返り，不明点があれば教科書や授業ノートで確認しておくこと。心理統計学の基礎を確実に身につけることは，今後の心理学研究において重要な基盤となるため，試験を通じて得られた知識の定着を図ることが重要である。また，試験の結果を踏まえて，自分の弱点や苦手分野を特定し，今後の学習計画に活かすことも有益である。


\addcontentsline{toc}{section}{\bibname}

% \bibliographystyle{../jecon_jpa}
% \bibliography{../sbib}
\printbibliography[title=引用文献]
\end{document}